% Article scientifique en version de travail, mis à jour le 29 septembre 2026.
% Compilation conseillée : latexmk -pdf -interaction=nonstopmode -halt-on-error
%                         -outdir=/tmp/garamon-article-tex Garamon_article_recherche_2026-09-27.tex
\ifdefined\pdfminorversion\pdfminorversion=7\fi
\documentclass[11pt,a4paper]{article}
\usepackage[T1]{fontenc}
\usepackage[utf8]{inputenc}
\usepackage[french]{babel}
\usepackage{lmodern}
\usepackage{amsmath,amssymb}
\usepackage{booktabs,tabularx}
\usepackage{graphicx}
\usepackage{pgfplots}
\usepackage{placeins}
\usepackage{needspace}
\pgfplotsset{compat=1.18}
\usepackage[margin=22mm]{geometry}
\usepackage[colorlinks=true,urlcolor=blue,linkcolor=blue,citecolor=blue]{hyperref}
\Urlmuskip=0mu plus 1mu\relax
\newcommand{\Cl}{\mathrm{Cl}}
\newcommand{\supp}{\operatorname{supp}}
\newcommand{\code}[1]{\texttt{#1}}
\newcommand{\figureattente}[4]{%
\begin{figure}[ht]
\centering
\fbox{\begin{minipage}[c][30mm][c]{.84\linewidth}
\centering\small\textbf{#1}\\[4pt]
\emph{Données à acquérir}\\[3pt]#2
\end{minipage}}
\caption{#3 Aucun point ni classement n'est représenté dans cette version.}
\label{#4}
\end{figure}}

\title{Calcul exact adaptatif dans Garamon.jl :\\
audit du générateur C++ et premières évaluations}
\author{}
\date{29 septembre 2026}

\begin{document}
\maketitle
\begin{abstract}
Le coût d'un calcul dans une algèbre de Clifford dépend de la dimension
ambiante, mais aussi des lames effectivement présentes, des coefficients
demandés et du nombre de réutilisations. Nous auditons le générateur C++
Garamon et étudions une réimplémentation Julia pure qui choisit entre
produit complet, parcours creux, jointure de supports, plans préparés et
espaces de travail. Nous définissons des scénarios de demandes de calcul,
un contrat de correction et des coûts séparés de préparation, compilation
et exécution. Un balayage de 3\,256 cas exacts couvre des dimensions de 2
à 256 ; le classement entre récursion et jointure ciblée s'inverse lorsque
le nombre de coefficients demandés change. Sur 192 cas de produit ternaire,
un plan préparé perd au premier appel mais gagne les vingt cellules
évaluées à 1\,024 appels. Une étude de workspace valide 687 cas, dont
150 noyaux réchauffés sans allocation observée. Ces temps ont été
recueillis en présence de calculs concurrents et restent exploratoires.
Les prototypes de rang binaire, radical nilpotent, Pfaffien et identités
polynomiales montrent des domaines d'admission distincts, sans établir
de seuil de vitesse général. La comparaison C++/Julia strictement
appariée vérifie la correction d'un noyau à tableaux compacts ; elle ne permet pas
encore de conclure à une supériorité de langage. Le prévol reproductible
de 32 voies Julia et de deux voies C++/Julia est prêt pour une campagne
isolée, qui devra confirmer les classements et leurs coûts complets.
\par\smallskip\noindent\textbf{Mots clés :} algèbre géométrique,
produit de Clifford, calcul creux, jointure, workspace, compilation Julia.
\end{abstract}

\section{Introduction}
Une algèbre $\Cl(V,G)$ de dimension $n$ possède $2^n$ coordonnées de base.
Cette taille rend une représentation dense coûteuse dès que $n$ croît ; elle
ne dit pourtant pas combien de coordonnées interviennent dans une requête.
Deux produits dans la même algèbre peuvent demander des calculs très
différents : un coefficient scalaire d'entrées creuses, une sortie complète
de blocs homogènes ou une longue série sur des supports inchangés. Une
optimisation valable pour l'une de ces demandes peut faire perdre du temps
ou de la mémoire à une autre. La question centrale est donc le choix
d'une méthode \emph{en fonction de la demande de calcul}, et non d'un seul
seuil dimensionnel.

Nous fixons trois exigences. Premièrement, une voie dite exacte doit
accumuler chaque contribution à chaque coefficient demandé ; avec des
coefficients flottants, cette exigence algorithmique ne supprime pas
l'arrondi. Deuxièmement, toute comparaison inclut le contrat de sortie,
la précision, la préparation, la compilation et la propriété des tampons.
Troisièmement, les résultats recueillis sous charge concurrente servent à
formuler des hypothèses, tandis que les seuils de sélection attendent un
rejeu isolé. L'élagage probabiliste est étudié séparément comme calcul
approché avec erreur observable et éventuelle restauration.

La contribution de cet article est une lecture du code C++ à une révision
fixée, une famille d'algorithmes Julia avec leurs domaines d'admission, puis
un protocole d'évaluation commun. Après les notions algébriques et la
cartographie des scénarios, nous distinguons les chemins disponibles en
C++ et Julia. Nous présentons ensuite les méthodes implémentées, les
transpositions encore expérimentales, les mesures obtenues et ce qu'elles
ne démontrent pas. Les références scientifiques expliquent l'origine des
idées ; chaque chiffre rapporté est défini ici par son scénario et son
contrat, sans supposer la lecture d'un document de travail annexe.

Les graphiques chiffrés de cette version sont des mesures exploratoires
et leurs légendes précisent les conditions de recueil. Les cadres marqués
« Données à acquérir » réservent la place des comparaisons encore ouvertes :
ils ne contiennent ni points simulés ni performance présumée. Le graphique
des bibliothèques GA sera inséré automatiquement quand la campagne
DrWatson aura produit son PDF vectoriel.

\section{Travaux antérieurs}
Le générateur Garamon combine calculs préétablis et parcours récursif sur un
arbre préfixe virtuel~\cite{breuils2019,garamoncpp}. Les bornes arithmétiques
de Breuils, Nozick et Sugimoto concernent deux multivecteurs homogènes pleins
\cite{breuils2020} ; leur nombre d'opérations ne prédit pas directement le
coût d'entrées creuses ou d'une seule sortie demandée. Gaalop compile des
expressions d'algèbre géométrique dans un autre cadre de spécialisation
\cite{gaalop}. Notre point de départ est plus large que la génération
d'une algèbre fixe : nous cherchons aussi à choisir un chemin à l'exécution
et à amortir sa préparation sur des séquences de demandes.

Les jointures de bases de données~\cite{leapfrog} motivent la recherche
d'un partenaire de masque déterminé par XOR, plutôt que l'énumération de
tous les produits intermédiaires. Les compilateurs de tenseurs creux
\cite{taco,finch,galley} et les workspaces modulaires
\cite{sparseworkspace} montrent que l'organisation des écritures et de la
sortie peut être aussi décisive que le nombre d'opérations arithmétiques.
Ces analogies guident nos prototypes ; la jointure de Clifford doit encore
traiter signes, grades, métrique et sorties possédées, et elle n'hérite
d'aucune borne de ces autres domaines sans nouvelle preuve. Les trains
tensoriels~\cite{oseledets} et diagrammes de décision à zéros supprimés
\cite{minato} suggèrent des représentations compactes seulement lorsque
les coefficients possèdent la structure requise.

La planification de FFTW~\cite{fftw} et les portefeuilles de solveurs
\cite{satzilla,autofolio} motivent un sélecteur dont le coût de décision
doit être inclus. L'égalité saturée~\cite{egg} propose une manière de
chercher des expressions équivalentes sous des règles prouvées ; SHAP
\cite{shap} peut expliquer un modèle de sélection après apprentissage,
sans certifier une décision. Le calcul approché et le replay ont un autre
contrat que le produit exact. Aucune performance de ces travaux n'est
attribuée à Garamon par simple analogie : les scénarios et oracles de la
section~\ref{sec:scenarios} rendent explicite ce qui doit être revalidé.

\section{Méthode : produits exacts et sélection de chemins}
Soit $V$ un espace vectoriel réel de dimension $n$ muni d'une forme
bilinéaire symétrique de matrice de Gram $G$. Son algèbre de Clifford est
définie par $uv+vu=2u^{\mathsf T}Gv$. Pour un sous-ensemble ordonné
$A=\{i_1<\cdots<i_k\}$, la lame $e_A=e_{i_1}\wedge\cdots\wedge e_{i_k}$
a le grade $k$ ; $e_{\varnothing}=1$ est le scalaire. Un multivecteur
$x=\sum_{A\subseteq\{1,\ldots,n\}}x_Ae_A$ possède au plus $2^n$
coefficients, mais son support $\supp(x)=\{A:x_A\ne0\}$ peut être bien
plus petit. Un bloc de grade $k$ contient au plus $\binom nk$ coefficients.
Le \emph{produit géométrique} prolonge la règle de Clifford ; le produit
extérieur en retient le grade $|A|+|B|$ lorsqu'il existe. Une requête
peut demander le produit entier, un ensemble de grades ou seulement les
coefficients d'un ensemble $O$ de masques. Dans ce dernier cas, la sortie
mathématique est la projection $P_O(xy)$ : les coefficients hors de $O$
ne font pas partie de la requête.

Pour une base orthogonale $e_1,\ldots,e_n$ avec $e_i^2=g_i$, on note $e_A$ la
lame ordonnée de masque $A$. Alors
\begin{equation}
e_Ae_B=\sigma(A,B)\left(\prod_{i\in A\cap B}g_i\right)e_{A\triangle B},
\qquad \sigma(A,B)\in\{-1,1\},
\label{eq:blade}
\end{equation}
où $\triangle$ est la différence symétrique. Un zéro de métrique, un grade de
sortie impossible ou une sortie non demandée permettent de supprimer un chemin
\emph{par preuve}. Pour le coefficient $K$ de $(ab)c$, les supports candidats
doivent vérifier $A\triangle B\triangle C=K$. En choisissant les deux plus
petits supports à parcourir, le troisième masque est déterminé par XOR ; les
facteurs et signes restent toutefois évalués dans l'ordre original du produit.
Cette jointure \code{join3} n'échantillonne aucun triplet admissible. Elle
s'inspire des parcours de jointure relationnelle~\cite{leapfrog}, sans être une
implémentation de Leapfrog Triejoin ni hériter de ses bornes.

Les exemples \emph{EGA} ci-dessous utilisent une métrique euclidienne
$G=I$. Une signature \emph{mixte} comporte des directions de carré
positif et négatif ; une signature \emph{dégénérée} possède au moins une
direction nulle. Ces trois cas ne sont pas interchangeables : un facteur
$g_i=0$ annule une contribution de~\eqref{eq:blade}, et une métrique
non diagonale introduit des contractions croisées. Les coûts de conversion
de base sont donc mesurés ou déclarés séparément.

Deux réalisations doivent être distinguées. La stratégie d'expression
\code{strategy=:join3} recalcule une sortie ciblée et a été chronométrée
ci-dessous. Le nouveau \code{TripleJoinPlan} prépare une fois tous les chemins
de plusieurs sorties, et \code{TripleJoinWorkspace} réemploie les tableaux de
coefficients sous budget de sondes, de chemins et de mémoire. Son résultat est
emprunté au workspace et écrasé à l'appel suivant. Les tests
\code{test/triplejoin.jl} passent 223 assertions sur 2, 3, 4, 6, 8, 12,
64, 65, 66, 127, 128 et 129 dimensions, y compris métrique dégénérée,
validations de supports et refus non orthogonal. Une première comparaison
PerfChecker de 192 cas lui est désormais propre ; les chiffres de
\code{join3} dans le balayage dimensionnel initial ne lui sont pas attribués.

La relation~\eqref{eq:blade} ne s'applique pas telle quelle à une base non
orthogonale : $e_i e_j=G_{ij}+e_i\wedge e_j$ pour $i\ne j$. Il faut alors un
produit direct adapté ou un changement de base dont les coûts et l'erreur
numérique sont inclus. La propagation conservatrice des grades à travers une
expression peut éviter l'énumération de supports intermédiaires impossibles et
faire admettre certains cas sous un budget de supports, mais elle ajoute son
propre coût de préparation. Les travaux sur les produits homogènes
\emph{pleins}~\cite{breuils2020} donnent des bornes pour un autre régime de
charge : on ne les transfère pas sans preuve à quelques lames dispersées.

\section{Scénarios de demandes de calcul}
\label{sec:scenarios}
\subsection{Ce qui constitue un scénario}
Un scénario est une \emph{suite de demandes algébriques} et non une
application géométrique nommée. Il fixe la métrique, les types numériques,
les supports et leurs changements, les opérations, les sorties requises,
l'horizon $H$ et les ressources. Ainsi, « produit de deux vecteurs » ne
détermine pas encore le coût : on peut demander le scalaire seulement,
le bivecteur, le multivecteur complet ou une chaîne de produits sur les
mêmes vecteurs. De même, une dimension $n=65$ avec deux lames actives
reste très différente d'un grade homogène plein de dimension moindre.

Nous distinguons l'horizon \emph{à froid} (processus neuf, chargement,
construction et première compilation), le premier appel dans un processus
déjà chargé, et un épisode \emph{chaud} de $H$ appels. Une charge est dite
\emph{stable} lorsque métrique et supports se répètent, même si les
coefficients changent ; elle est \emph{variable} lorsque supports, grades
ou métrique changent et invalident potentiellement un plan. La sortie
\emph{possédée} survit à l'appel suivant ; une vue empruntée à un workspace
doit être consommée immédiatement ou copiée. Ces distinctions changent le
coût total et sont conservées dans les comparaisons.

\begin{table}[ht]
\centering\small
\begin{tabularx}{\linewidth}{@{}p{30mm}X X@{}}
\toprule
Demande de calcul & Variables qui changent le travail & Voies à confronter \\
\midrule
Produit complet ponctuel & Faible $n$, entrées denses ou quelques lames, sortie possédée, $H=1$. & Direct dense, creux, récursif et génération courte ; compter préparation et première compilation. \\
Produit homogène & Grades d'entrée fixés, de presque vide à plein ; sortie complète ou grades choisis. & Blocs de grades et parcours direct, avec coût de conversion. \\
Projection ciblée & Un scalaire, un coefficient $K$ ou quelques sorties de $ab$ ou $(ab)c$. & Jointure XOR ou ternaire, récursion filtrée et extraction du produit complet ; vérifier chaque contribution admise. \\
Épisode stable & Même métrique et mêmes supports ; coefficients renouvelés sur $H=32$ ou $1\,024$ appels. & Plan préparé, workspace, noyau généré et cache ; inclure construction et possession des sorties. \\
Trace variable & Supports, grades et parfois métrique changent entre appels. & Recalcul direct et cache de plans LRU ou roulette ; mesurer ratés, évictions et mémoire retenue. \\
Processus neuf & Algèbre et support connus à l'avance, ou découverts après le chargement. & JIT à la demande et catalogue natif fini ; séparer chargement, reconstruction et compilation. \\
Structure et grande dimension & Peu de directions actives, lames factorisées, chaîne scalaire ou sortie qui ne tient pas en RAM. & Rang binaire, versor, Pfaffien, TT/ZDD exacts, flux par blocs ou refus de budget. \\
Exécution parallèle & Beaucoup de produits indépendants, lots compacts, CPU multi-thread ou GPU. & Un thread, tâches, processus et CUDA ; inclure transfert, synchronisation et duplication des tampons. \\
Métrique générale & Signature mixte, dégénérée ou base non orthogonale ; même sortie contrôlée. & Produit direct, transformation extérieure, quotient nilpotent si admissible. \\
Calcul approché explicite & Un utilisateur accepte une erreur sur une récurrence ou un coefficient. & Élagage probabiliste et replay face au calcul exact ; publier erreur et coût de restauration. \\
\bottomrule
\end{tabularx}
\caption{Cartographie des demandes utilisées pour interpréter les essais. Elle
définit des comparaisons, pas une victoire présumée d'une méthode. Les
stratégies structurelles exigent leurs certificats d'admission.}
\label{tab:scenarios}
\end{table}

Une comparaison \emph{appariée} réutilise, pour chaque cellule, la même
algèbre, la même séquence d'entrées produite avec une graine fixée, le
même type numérique et la même sortie mathématique. Elle signale si la
sortie est copiée, si un cache est déjà chaud et si la préparation est
incluse. Les stratégies stochastiques reçoivent également une graine
fixée et leurs paramètres sont enregistrés ; cette graine rend le rejeu
possible, sans rendre exact un élagage approché. Une limite de RAM, de
VRAM ou de temps mène à un \emph{refus documenté}, jamais à l'omission
silencieuse d'un terme exact. Les comparaisons par dimension gardent
séparés les scénarios du tableau~\ref{tab:scenarios} : une victoire pour
une sortie scalaire ne vaut pas pour la sortie complète.

\section{Comparatif documenté de Garamon C++ et de Garamon.jl}
Le code C++ étudié est le commit public \code{6267161} du dépôt
Garamon~\cite{garamoncpp}, en regard de l'article du générateur
\cite{breuils2019}. Le générateur lit une configuration et une métrique, prépare
l'indexation masque--(grade,position), diagonalise si nécessaire, construit
des matrices creuses de changement de base par grade et émet une bibliothèque
C++. Son multivecteur conserve les grades présents dans une liste, avec un
vecteur de coefficients dense pour chaque grade. Les produits choisissent des
fonctions déroulées ou des parcours récursifs d'un arbre préfixe \emph{virtuel}
suivant le grade et le seuil de génération. Les chemins lus ne matérialisent
pas un arbre de nœuds de coefficients et ne mémoïsent pas les valeurs des
sous-produits~\cite{cppmain,cppindex,cppdispatch,cppmvec,cppgeometric,cppmetric}.

Le code Julia actuel, publié sur la branche \code{bench-ready-20260929}
de \code{Garamon.jl}, est une
réimplémentation Julia pure, sans wrapper C++. Il dispose de multivecteurs
denses bornés et de dictionnaires de coefficients actifs, de produits directs,
de plans \code{ProductPlan}, d'un DSL d'expressions, de sorties ciblées et de
génération Julia bornée. Les masques sont actuellement \code{UInt64} jusqu'à
64 directions, \code{UInt128} de 65 à 128, puis \code{BigInt}. Le parcours
orthogonal creux n'exige pas de table globale $2^n$. Le stockage dense explicite
est limité à $n\leq20$ ; cela n'implique pas que tous ces cas soient rapides ou
admis par les budgets de produit.

Le code Julia vérifié comprend \path{src/algebra.jl} (métriques),
\path{src/multivectors.jl} (stockages et masques) et \path{src/plans.jl}
(plans). Les coefficients ciblés et les expressions relèvent de
\path{src/products.jl} et \path{src/dsl.jl}. La génération et les tampons
réutilisables relèvent de \path{src/generated.jl}, \path{src/workspace.jl} et
\path{src/triplejoin.jl}. Ces fichiers appartiennent à l'arbre de travail
Julia à la révision \code{1a4cb87}, au-delà du commit initial
\code{adf1dc5}.

\begin{table}[ht]
\centering\small
\begin{tabularx}{\linewidth}{@{}p{31mm}X X@{}}
\toprule
Aspect & Garamon C++ lu & Garamon.jl observé \\
\midrule
Préparation & Génération puis compilation d'une bibliothèque pour une algèbre configurée. & Algèbre à l'exécution, plan de chemins sur demande ; spécialisation Julia facultative, au plus 256 chemins déroulés. \\
Coefficients & Grades présents, vecteur dense dans chaque grade. & Dense global borné ou seuls masques actifs ; blocs homogènes préparés possibles. \\
Grande dimension & Tables d'indexation $2^n$ et masques \code{unsigned int} dans la révision lue. & Masques de largeur variable et budgets explicites ; la faisabilité dépend des supports. \\
Base non\newline orthogonale & Diagonalisation puis matrices de transformation extérieure par grade. & Produit général direct et transformation extérieure à la demande, sous budget. \\
Réutilisation & Code généré par algèbre ; pas de cache de sous-produits observé dans les chemins audités. & Plan, cache de plans, programme généré et workspace distincts. \\
\bottomrule
\end{tabularx}
\caption{Comparaison de conception fondée sur les fichiers audités. Elle ne compare pas les vitesses des deux langages.}
\label{tab:cppjulia}
\end{table}

Les écarts de sémantique doivent rester explicites : un inverse singulier
déclenche une erreur dans le port Julia, tandis que la voie C++ auditée renvoie
zéro près de son seuil. Les résultats de performance du noyau public Julia et
ceux du banc C++/Julia ont des contrats distincts ; la couverture de tous
les scénarios C++ n'est pas déduite des seuls tests présentés. Le
tableau~\ref{tab:cppjulia} compare les capacités observées ; les sections
suivantes définissent leurs domaines et les essais correspondants.

Nous appelons \emph{packed} un noyau qui lit les chemins préparés et les
coefficients dans des tableaux contigus, puis écrit une matrice de sortie
possédée. Une première comparaison appariée a porté cette boucle Julia dans
un noyau C++ \emph{expérimental}. Elle garde mêmes chemins, ordre
d'accumulation, tableaux, coefficients, sorties et précision ; les mesures
chaudes passent par le même frontend Julia. Le vrai \code{Mvec<double>} issu
du générateur C++ est chronométré à titre \emph{descriptif}, avec son
stockage par grades, et ne forme pas un duel de langages avec la voie packed.
L'extension de cette comparaison attend d'abord la validation et
l'optimisation expérimentale de chaque technique Julia ; elle exigera
ensuite un noyau C++ de même stratégie et l'alignement de la préparation,
de la compilation et du chronométrage. Les bibliothèques
C++ ont été construites dimension par dimension dans un budget surveillé,
avec relevé de disque et nettoyage des artefacts de chaque dimension ;
la répétition isolée reste à faire.

\subsection{C++ et Julia : premier noyau packed apparié}
\label{sec:cppperf}
\paragraph{Pourquoi cet essai ?}
Comparer deux langages avec des algorithmes différents confond le coût
du langage et celui de la stratégie. Ce pilote porte donc le même noyau
packed et vérifie d'abord l'identité des entrées, sorties et opérations
avant d'examiner ses temps.
Une première campagne exécutable valide 64/64 cas PerfChecker, avec
31 échantillons chauds par cas, dans EGA3 et EGA7. Deux familles d'entrées
(\emph{mixed}, dont les supports mêlent plusieurs grades, et
\emph{vectors}, limitées au grade~1) et quatre tailles de lots indépendants
($H=1,32,1024,10000$) sont couvertes. L'oracle entier indépendant vérifie
chaque coefficient et la complétude du support. Il ne s'agit pas d'une
récurrence. Les trois modalités \emph{appariées} emploient la même boucle
de chemins, le même ordre des sommes, les mêmes matrices \code{Float64}
et la même allocation de sortie via le frontend Julia : la fonction publique
\code{run\_packed\_batch} en Julia JIT, la même fonction après
\code{Base.precompile} dans le processus, et un port C++ \emph{expérimental}
de ce noyau. Le C++ packed est compilé avec Clang 18.1.3,
\code{-O2}, \code{-march=native} et contraction flottante désactivée.
Ces choix bornent la comparaison à une même stratégie
et ne représentent pas le dispatch C++ amont par grades.

\begin{figure}[ht]
\centering
\begin{tikzpicture}
\begin{axis}[
  width=.91\linewidth,height=6.2cm,
  ybar,bar width=7pt,enlarge x limits=.15,
  ymin=0,ymax=.86,ytick={0,.2,.4,.6,.8},
  xtick={1,2,3,4},
  xticklabels={3D mixte,3D vect.,7D mixte,7D vect.},
  ylabel={Lot de 10\,000 produits (ms)},
  grid=major,grid style={gray!20},
  legend style={at={(.37,.98)},anchor=north west,legend columns=1,draw=none,fill=white},
]
\addplot+[fill=blue!55,draw=blue!80!black] coordinates
  {(1,.566444) (2,.138876) (3,.362696) (4,.741087)};
\addplot+[fill=gray!45,draw=black] coordinates
  {(1,.604970) (2,.142643) (3,.374018) (4,.756926)};
\addplot+[fill=orange!65,draw=orange!85!black] coordinates
  {(1,.543365) (2,.140069) (3,.358855) (4,.685747)};
\legend{Julia JIT,Julia précompilée,C++ packed}
\end{axis}
\end{tikzpicture}
\caption{Médianes chaudes pour quatre lots de 10\,000 produits, dans un même
contrat packed et via le même frontend. « Précompilée » signifie
\code{Base.precompile} dans le processus. Cette première passe était
concurrente avec Étendue3D, sans affinité CPU ; l'écart entre barres ne
constitue pas une victoire générale.}
\label{fig:cpppacked}
\end{figure}

Les temps restent proches dans ces quatre familles : par exemple, en EGA3
\emph{vectors}, les médianes Julia JIT / C++ packed sont
0,138876/0,140069~ms. Le \emph{vrai} produit
\code{Mvec<double>} d'une bibliothèque générée par Garamon a aussi été
exécuté et vérifié sur 16 charges : à $H=10000$, ses médianes sont
2,724, 0,751, 1,746 et 1,385~ms dans l'ordre de la
figure~\ref{fig:cpppacked}. Son stockage par grades, son dispatch et ses
allocations diffèrent ; aucun ratio C++/Julia n'est déduit de ces nombres.
L'exécutable C++ natif décompose encore le noyau packed, mais ses temps
ne sont pas soustraits aux temps du frontend pour fabriquer un « coût FFI ».
\figureattente{Comparaison packed isolée, C++ / Julia}
{Abscisse : dimension, famille de supports et horizon ; ordonnée : temps
d'épisode apparié, avec préparation et première compilation séparées.}
{Rejeu du même noyau sur la machine dédiée, après les mesures exploratoires
de la figure~\ref{fig:cpppacked}.}{fig:cpppacked-isole}
Dans la voie \code{Mvec} auditée, l'insertion de certains vecteurs Eigen
temporaires se fait par copie. Une expérience \emph{C++ contre C++}
indépendante a maintenant mesuré le déplacement de ces temporaires
avec oracle, compteurs d'allocations et deux ordres temporels
(section~\ref{sec:moveperf}). Elle ne modifie pas la comparaison packed
C++/Julia de la figure~\ref{fig:cpppacked}.

Le premier appel Julia JIT observé dure 31,68--34,30~ms ; presque tout
ce temps est attribué à la compilation. L'appel à
\code{Base.precompile} coûte 36,03--51,61~ms en amont. Les premiers
produits correspondants rapportent ensuite zéro JIT. Cela
ne démontre ni gain froid total, ni code conservé entre sessions. Le build du
générateur C++ prend environ 18,66~s ; la bibliothèque packed environ
0,20~s, et la bibliothèque amont 1,32~s en EGA3 ou 3,45~s en EGA7.
Les artefacts produits ont été mesurés (3\,628\,956 et
7\,065\,214 octets respectivement), puis supprimés avant la dimension
suivante. Le plafond surveillé est de 256~Mio par dimension. Ces phases ne sont pas additionnées
à une médiane chaude pour inférer un seuil.

L'instabilité de la première passe est visible en 7D \emph{mixed},
$H=1024$ : la \emph{même} fonction chaude Julia donne
17,452~\(\mu\)s après JIT ordinaire et 34,956~\(\mu\)s après
précompilation ciblée, contre 16,999~\(\mu\)s pour C++ packed.
Ce facteur proche de deux ne démontre aucun effet causal de la
précompilation. Les résumés des 31 échantillons sont conservés, mais pas
les temps individuels de cette passe ; ils ne peuvent être reconstruits.
La reprise prévue attend un CPU libre, fixe un cœur P,
alterne les ordres dans trois paires complètes et conserve les échantillons
individuels ainsi que l'activité système. Jusque-là, la figure
\ref{fig:cpppacked} est un \emph{criblage}, sans classement général.
\FloatBarrier

\subsection{C++ contre C++ : déplacer les buffers Eigen temporaires}
\label{sec:moveperf}
\paragraph{Pourquoi cet essai ?}
Le \code{Mvec} généré peut copier un buffer Eigen lorsqu'un vecteur
temporaire est inséré dans sa structure de grades. Le test isole ce
coût de copie dans le même langage et la même représentation, puis
vérifie que le déplacement conserve tous les produits.
Une modification isolée du \code{Mvec.hpp} amont remplace quatre insertions
d'un \code{Kvec} local par \code{std::move(kvec)} et inclut
\code{<utility>}. Deux sites concernent l'accès par indice ;
deux, la création paresseuse d'un vecteur de grade. Chaque temporaire n'est plus
utilisé après l'insertion : le déplacement transfère son buffer Eigen
au nœud de liste, au lieu de créer une seconde allocation et copie des
coefficients. L'ordre arithmétique, les supports, les grades et les tests de
zéro restent identiques. Le clone amont demeure intact ; la modification
est conservée dans un worktree C++ isolé, sans publication ni PR.

Le banc natif oppose \emph{le même C++} avant et après ce seul changement.
Il reprend EGA3/EGA7, supports \emph{mixed}/\emph{vectors},
$H=1,32,1024,10000$ et écrit dans une matrice préallouée tous les
coefficients structurellement possibles pour ces supports fixés.
Les entrées et l'oracle entier sont identiques dans les deux variantes ;
leur préparation et l'allocation de la matrice sont exclues du noyau.
Les exécutables amont et modifiés vérifient au total 67\,840 produits
exacts, y compris toutes les paires de lames de base et 256 paires de
multivecteurs supplémentaires par dimension et par exécutable ; les
tests inspectent tous les coefficients, y compris les zéros.
Un test séparé confirme le transfert d'adresse du buffer Eigen et
l'absence d'alias avec les opérandes. Le domaine démontré est
\code{Float64}, métrique euclidienne, 3D et 7D.

\begin{figure}[ht]
\centering
\begin{tikzpicture}
\begin{axis}[
  width=.84\linewidth,height=5.8cm,
  ybar,bar width=11pt,enlarge x limits=.23,
  ymin=0,ymax=100,ytick={0,20,40,60,80,100},
  xtick={1,2,3,4},
  xticklabels={3D mixte,3D vect.,7D mixte,7D vect.},
  ylabel={Appels malloc Eigen (milliers)},
  grid=major,grid style={gray!20},
  legend style={at={(.98,.98)},anchor=north east,draw=none,fill=white},
]
\addplot+[fill=blue!55,draw=blue!80!black] coordinates
  {(1,86.666) (2,40) (3,60) (4,40)};
\addplot+[fill=orange!65,draw=orange!85!black] coordinates
  {(1,43.333) (2,20) (3,30) (4,20)};
\legend{C++ amont,C++ avec déplacements}
\end{axis}
\end{tikzpicture}
\caption{Appels \code{malloc} des buffers Eigen pour 10\,000 produits :
deux ordres et trois répétitions donnent les mêmes comptes. Les
allocations de nœuds \code{new} restent inchangées ; les appels
\code{malloc} plus \code{new} diminuent d'un tiers dans les 32 groupes
appariés. Il s'agit de comptes cumulés instrumentés, pas de pics RSS.}
\label{fig:cppmovealloc}
\end{figure}

Les 32 groupes appariés (16 charges dans deux ordres), soit 192 observations
de compteurs, divisent exactement par deux le nombre et les octets demandés
des buffers Eigen ; les allocations de nœuds ne changent pas. Les octets
\emph{cumulés} \code{malloc+new} en EGA3 \emph{mixed} à $H=10000$
passent de 5\,839\,960 à 3\,786\,640 ; les autres familles et horizons
suivent le même mécanisme. Les compteurs excluent entrées, oracle, sortie
et appels internes à des bibliothèques partagées non interceptés.
Le binaire de temps n'instrumente pas les allocations, afin de ne pas
ralentir le noyau qu'il chronomètre.

\begin{table}[ht]
\centering\small
\begin{tabular}{@{}lrr@{}}
\toprule
Charge, $H=10000$ & Amont puis déplacement & Déplacement puis amont \\
\midrule
EGA3 mixte & 2,525 $\to$ 2,217 & 2,441 $\to$ 1,981 \\
EGA3 vecteurs & 0,692 $\to$ 0,527 & 0,689 $\to$ 0,521 \\
EGA7 mixte & 1,310 $\to$ 1,013 & 1,307 $\to$ 1,019 \\
EGA7 vecteurs & 0,951 $\to$ 0,807 & 0,933 $\to$ 0,784 \\
\bottomrule
\end{tabular}
\caption{Médianes de 101 observations par variante et ordre, millisecondes
par lot ; la flèche va toujours de l'amont au C++ modifié. Temps
exploratoires sous activité Étendue3D, distincts de la comparaison
C++/Julia packed.}
\label{tab:cppmovetimes}
\end{table}

Les temps favorisent le déplacement dans ces deux passages, mais la
concurrence Étendue3D subsistait malgré l'affinité du contrôleur au
cœur P logique 0. Le cas EGA7 vecteurs, $H=32$, donne par exemple des
ratios modifié/amont de 0,985 puis 0,827 selon l'ordre. Les 101 temps
individuels sont conservés pour cette nouvelle campagne, contrairement au
premier banc packed. Sans contrôle d'activité complet ni calibration
avant/après, seule la réduction d'allocations du domaine testé constitue
ici une conclusion stable ; l'amplitude temporelle attend une reprise
isolée. Les objets générés de chaque dimension ont été supprimés avant la
suivante, sous le même plafond disque surveillé de 256~Mio.
\FloatBarrier

\section{Plans, workspaces, JIT et durée de vie}
Un plan structurel garde les masques, chemins $(i,j,k)$ et facteurs métriques,
pas les valeurs des opérandes. Un \code{ProductWorkspace} préalloue deux vecteurs
d'entrée, un vecteur de sortie et un multivecteur creux de résultat. Chaque
appel revalide l'algèbre et le support, rafraîchit les coefficients, remet les
sorties à zéro et accumule tous les chemins. Le résultat retourné est emprunté
au workspace : l'appel suivant l'écrase ; garder un historique exige une copie.
Une instance mutable appartient à une tâche concurrente. L'admission compare
la taille atteignable de l'objet à un \code{max\_bytes} et à un quart de la
mémoire système libre observée à la construction. Il s'agit d'un contrôle
ponctuel, ni d'un plafond de RAM du processus, ni d'une borne sur le code JIT
ou la croissance des coefficients à précision arbitraire. Le travail sur les
workspaces creux~\cite{sparseworkspace} motive cette expérimentation, sans
garantir le meilleur format pour Garamon.

Quatre durées de vie sont distinctes : (1) la description de l'algèbre et son
éventuel changement de base, (2) le plan structurel, (3) les tampons mutables,
(4) le code machine compilé. Le produit généré encode ses indices de chemins
dans un paramètre de type Julia, tandis que facteurs et coefficients demeurent
des valeurs ; préparer l'objet n'équivaut pas à compiler sa première exécution.
La précompilation du package ou une image applicative peuvent couvrir des
spécialisations représentatives, mais un nouveau support ou un nouveau type
peut encore déclencher du JIT~\cite{juliaprecompile}. La politique doit être
évaluée pour une algèbre connue avant démarrage et pour des structures
découvertes en cours de session. L'étude du catalogue natif ci-dessous examine ces
deux situations ; la sauvegarde d'un \code{ProductPlan} ne prouverait pas à
elle seule la persistance du code natif entre processus.

Pour $H$ appels dans un même processus et à sortie identique, on compare
\begin{equation}
C_{\mathrm{plan}}+C_{\mathrm{workspace}}+C_{\mathrm{JIT,prep}}
  +H\,t_{\mathrm{prep}}
\quad\text{à}\quad C_{\mathrm{JIT,direct}}+H\,t_{\mathrm{direct}},
\label{eq:amort}
\end{equation}
en ajoutant aux deux côtés chargement, conversion, copie et conservation des
résultats lorsque le scénario les exige. Si la compilation est déjà payée dans
le scénario chaud, les deux termes JIT s'annulent ; dans un processus neuf,
ils doivent être mesurés séparément. Il s'agit d'un modèle de décision, pas
d'un seuil déduit des campagnes ci-dessous.
L'approche par portefeuille, comme les plans de FFTW~\cite{fftw}, n'est
pertinente que si son coût de décision et de préparation s'amortit sur les
charges réelles ; elle ne garantit aucun gain par analogie.

\section{Protocole expérimental et portée des résultats}
\label{sec:protocol}
Les comparaisons de cet article suivent le scénario du
tableau~\ref{tab:scenarios}. Pour chaque cellule, une graine produit les
entrées et la trace de demandes avant l'appel des stratégies. Un oracle
indépendant énumère les petits mots de Clifford avec entiers bornés,
compte la parité des inversions de directions et les facteurs métriques,
puis confronte tous les coefficients \emph{demandés}. Dans les fixtures
\code{Float64} choisies pour cet oracle, les sommes entières restent
représentables exactement ; ce fait local ne certifie pas tout calcul
flottant. Les métriques non diagonales, inverses, changements de base
et sorties structurées reçoivent en outre leurs identités propres ou
des refus de domaine contrôlés. Les tests de correction vérifient la
sortie possédée quand le scénario l'exige, ainsi que l'invalidation
des plans et la reprise d'une campagne interrompue.

Nous séparons quatre niveaux de preuve. Une identité algébrique démontre
qu'un chemin exact ne perd aucun terme dans son domaine. Un prévol
vérifie cette identité sur un corpus borné et des oracles indépendants ;
son compte de cas n'est pas une mesure de vitesse. Un profil instrumenté
localise des allocations, instabilités de type, coûts CPU ou GPU, mais
peut perturber le temps du noyau. Enfin, une comparaison de vitesse
demande des échantillons appariés, des ordres alternés, la même sortie et
une machine isolée. La dernière étape reste à conduire pour les grilles
étendues ; les médianes déjà présentées sont des \emph{écrans
exploratoires}, pas des seuils de sélection publiables.

Le banc \code{GaramonBench.jl} emploie Julia~1.13 et
DrWatson.jl~\cite{drwatson} pour produire des identifiants de cas et
archiver configuration, graine, environnement, sources, exclusions de
budget, oracles et échantillons. Une relance vérifie l'intégrité d'un
cas achevé avant de le sauter ; un changement de source ou de machine
requiert une nouvelle racine de résultats. Dans l'état source étudié,
un prévol minimal raccorde 32 voies Julia (30 sur un thread, une sur
quatre threads, une sur GPU) et deux cas packed C++/Julia à contrat
identique. Les 32 profils PerfChecker correspondants sont complets et
audités ; quatre cas GPU possèdent chacun des traces distinctes de
noyau résident, de retour hôte et d'épisode avec transferts. Cette
qualification de l'instrumentation ne remplace ni le balayage exhaustif
des paramètres ni la répétition isolée. Les prototypes structurels
décrits plus loin peuvent avoir leurs propres corpus et budgets : leurs
comptes ne sont pas additionnés aux 32 cas du prévol minimal.

La machine des premiers essais est un Intel Core i7-12700 avec une
RTX~3060~Ti de 8~Gio. D'autres calculs ont parfois tourné en parallèle.
Nous rapportons donc les temps avec leur dimension, horizon, sortie et
nombre d'observations lorsque ces données existent, et nous n'agrégeons
pas les microsecondes de campagnes distinctes. Pour une comparaison
de langages, seule la voie packed à mêmes chemins et sorties est
appariée ; le \code{Mvec} du générateur C++ est une autre stratégie.
Cette séparation évite d'attribuer au langage une différence de
représentation ou de génération.

\section{Techniques implémentées : question, essai et résultat}
\label{sec:techniques}
Chaque technique est présentée avec le problème qui motive son choix,
puis avec l'expérience correspondante, ses chiffres et sa limite de
validité. Les figures rapprochent des routes qui calculent la même
sortie ; elles ne comparent pas les durées de campagnes différentes.
Dans cette section, \emph{implémenté} signifie présent dans l'arbre Julia
étudié, \emph{mesuré} signifie évalué sur un corpus et un contrat de sortie
nommés, et \emph{hypothèse} signifie encore sans vitesse démontrée. Les
références primaires motivent les mécanismes ; les chiffres restent ceux de
Garamon. Un gain sur un coefficient ciblé n'est pas un gain sur le
multivecteur entier, et un temps chaud ne couvre pas automatiquement
préparation, compilation ou copie de résultat.

\subsection{Densité, supports creux et blocs de grades}
\paragraph{Pourquoi cette technique ?}
Le générateur C++ de Breuils, Nozick et Fuchs~\cite{breuils2019} conserve
les grades présents avec un tableau dense de taille $\binom nk$ par grade ;
Garamon.jl offre aussi un tableau dense global borné par $n\leq20$, un
dictionnaire des seuls masques actifs, des plans de blocs homogènes et une
liste triée pour certains produits extérieurs. Les compilateurs de tenseurs
creux TACO~\cite{taco} et Finch~\cite{finch} montrent l'intérêt de dissocier
format, parcours et calcul ; ils ne résolvent pas à eux seuls les signes,
annulations et changements de grade du produit de Clifford. Le format dense
vise les petits $n$ ou grades bien occupés, le dictionnaire les supports
dispersés, et le bloc de grade les appels répétés à structure homogène.
\paragraph{Essais, résultats et limites.}
Le dense paie des zéros et potentiellement $2^n$ cases ; le creux paie
recherche, hachage et localité moins régulière ; un bloc préparé paie
énumération et mémoire avant la première valeur. Sur une famille homogène en
12D, grade 3 par grade 3, le plan couvre 48\,400 chemins, occupe 1,81~Mo
et coûte 9,23~ms à préparer, puis 0,24~ms par produit entier ; huit supports
variables coûtent 11,02~ms préparation incluse, contre 27,14~ms en direct
et 19,36~ms avec reconstruction de plans propres à chaque support
. Pour un coefficient wedge sommital en 12D, la liste triée
\code{UInt64} prend 0,93 contre 5,42~\(\mu\)s sur 56 masques ; son prototype
\code{BigInt} en 66D perd face à la jointure. Ces constats
ne justifient pas un seuil universel de densité. Le protocole suivant doit
apparier même opération, sortie, supports et types, puis compter conversion,
construction du bloc, mémoire maximale et horizon réel ; il doit inclure des
grades centraux pleins et des supports changeants.

\subsection{Récursion préfixe C++ et planification des grades}
\paragraph{Pourquoi cette technique ?}
Dans le C++ audité, les produits sont déroulés pour certains petits grades
et parcourent sinon un arbre préfixe virtuel ; les nœuds ne contiennent pas
de coefficients persistants~\cite{cppgeometric,cppdispatch}. Ce choix évite
d'émettre toutes les combinaisons dans le code source. Les bornes de
Breuils, Nozick et Sugimoto~\cite{breuils2020} visent deux multivecteurs
homogènes \emph{pleins}. La récursion est donc une référence pertinente
pour les produits complets ou à nombreux grades, et la propagation de grades
peut éliminer des sorties impossibles avant un parcours de supports.
\paragraph{Essais, résultats et limites.}
Le branchement récursif, les indices par grade et les tableaux denses ont
un coût même lorsque peu de coefficients existent. Dans le balayage Julia
initial, en 256D dispersée, quatre sorties demandées prennent 140,7~\(\mu\)s
par récursion contre 223,6 par \code{join3}, mais pour la seule sortie
scalaire les temps sont 130,3 et 87,6~\(\mu\)s
(tableau~\ref{tab:dimensions}). C'est une comparaison entre chemins
\emph{Julia}, pas une mesure de la récursion C++ amont. Le premier banc
exécutable C++ mesure \code{Mvec} séparément et un port \emph{packed}
apparié à Julia, sans isoler cette récursion comme stratégie identique.
Un futur essai entre langages doit d'abord assurer identité de
stratégie, métrique, supports, coefficients, précision, sorties et contrat
de chronométrage, puis publier génération, compilation, première exécution,
exécution chaude et espace disque par dimension. La propagation des grades
reste à confronter à son propre coût de préparation, surtout si la demande
de sortie change.

\subsection{Jointure XOR et \code{TripleJoinPlan}}
\paragraph{Pourquoi cette technique ?}
Sous métrique diagonale, le masque de sortie d'un produit binaire est
$K=A\triangle B$ : pour $A$ et $K$ fixés, $B$ est déterminé. Dans
$(ab)c$, $C=A\triangle B\triangle K$. Le parcours de deux supports et la
sonde du troisième remplacent l'expansion des intermédiaires par une jointure
exacte. Cette idée est voisine des jointures à plusieurs relations
\cite{leapfrog}, sans en reprendre une borne de complexité. La cible est
une ou quelques sorties demandées, avec supports creux et algèbre diagonale.
Le \code{TripleJoinPlan} implémenté mémorise les chemins structurels
$(i,j,k)$ de plusieurs sorties ; son workspace rafraîchit les valeurs sans
partager les tampons entre tâches.
\paragraph{Essais, résultats et limites.}
Une jointure paie sondes, signes et facteurs métriques ; multiplier les sorties
peut refaire des parcours. L'expression \code{join3} mesurée bat le complet
sur le scalaire dispersé 256D (87,6 contre 1\,643,7~\(\mu\)s), mais perd
face à la récursion pour quatre demandes dans la même famille
(223,6 contre 140,7~\(\mu\)s). Les 223 assertions du
\code{TripleJoinPlan} couvrent sa correction sur les cas testés. Un premier
banc apparié de 192 cas compare maintenant le plan ternaire, \code{join3}
et deux produits complets sur les \emph{mêmes} trois opérandes et sorties :
avec préparation incluse, le plan perd pour $H=1$ dans les dix dimensions,
mais gagne les vingt cellules à $H=1024$
(section~\ref{sec:tripleperf}). Les supports n'ont que deux à quatre
directions actives et les échantillons longs sont parfois uniques.
La confirmation doit ajouter les deux plans binaires, la récursion, des
supports plus riches, la matérialisation éventuelle de chaque sortie et
une répétition isolée ; elle compte chemins admis, JIT et mémoire.
Les métriques non orthogonales exigent une autre formule et sont refusées
par cette jointure.

\subsection{Dimensions et sorties demandées}
\paragraph{Pourquoi cet essai ?}
La dimension seule ne dit pas si le calcul utile est dense ou creux.
Ce balayage demande quand une sortie ciblée évite assez de travail
pour compenser son adressage, et quand le produit complet reste préférable,
à supports et budgets explicités.
Le balayage continu a validé 3\,256 cas et consigné 6\,944 exclusions de budget.
Il couvre chaque entier de 2 à 256 avec au plus douze bivecteurs dispersés
plus le scalaire, 2 à 70 dans deux autres familles creuses, 2 à 11 pour le grade
2 plein et 2 à 5 pour l'entrée totalement occupée. Quatre stratégies et une ou
quatre sorties sont comparées. La limite 256 est une enveloppe expérimentale,
pas une saturation intrinsèque des produits creux. Les temps du
tableau~\ref{tab:dimensions} sont des médianes par lot, sans préparation ; la
première compilation et les allocations ne sont pas incluses dans ces
médianes et ne sont donc pas déduites du tableau.

\begin{table}[ht]
\centering\small
\begin{tabular}{@{}lrrrr@{}}
\toprule
Famille, $n$, sortie & Complet & Récursif & Grades & \code{join3} \\
\midrule
Creux dispersé, 2, scalaire & 1,9 & 9,8 & 14,6 & 5,2 \\
Creux dispersé, 5, scalaire & 29,3 & 33,3 & 40,9 & 17,6 \\
Creux dispersé, 64, scalaire & 211,5 & 34,4 & 41,4 & 20,5 \\
Creux dispersé, 65, scalaire & 1\,380,5 & 77,8 & 103,9 & 75,0 \\
Creux dispersé, 256, scalaire & 1\,643,7 & 130,3 & 130,3 & 87,6 \\
Creux dispersé, 256, quatre & 1\,674,3 & 140,7 & 144,8 & 223,6 \\
Grade 2 plein, 11, scalaire & 4\,915,9 & 277,8 & 339,8 & 269,7 \\
Grade 2 plein, 11, quatre & 4\,849,5 & 342,1 & 407,7 & 958,7 \\
\bottomrule
\end{tabular}
\caption{Microsecondes par lot dans le balayage dimensionnel initial. À
cette révision des sources, la dimension 65 passait de \code{UInt64} à
\code{BigInt} ; une campagne ultérieure a mesuré le chemin \code{UInt128}.}
\label{tab:dimensions}
\end{table}

\begin{figure}[ht]
\centering
\begin{tikzpicture}
\begin{axis}[
  width=.78\linewidth,height=5.6cm,
  ybar,bar width=14pt,enlarge x limits=.32,
  ymin=0,ymax=260,ytick={0,50,100,150,200,250},
  xtick={1,4},xticklabels={Une sortie,Quatre sorties},
  ylabel={Temps par lot (\(\mu\)s)},
  grid=major,grid style={gray!20},
  legend style={at={(.02,.98)},anchor=north west,draw=none,fill=white},
]
\addplot+[fill=blue!55,draw=blue!80!black] coordinates {(1,130.3) (4,140.7)};
\addplot+[fill=orange!65,draw=orange!85!black] coordinates {(1,87.6) (4,223.6)};
\legend{Récursion,\code{join3}}
\end{axis}
\end{tikzpicture}
\caption{En 256D dispersée, le nombre de sorties inverse le classement des
deux chemins Julia. Deux des quatre sorties sont nulles par structure dans
cette famille. Temps exploratoires du balayage initial, même révision à masques
\code{BigInt} pour ces deux cellules.}
\label{fig:sorties}
\end{figure}

Pour une seule sortie, \code{join3} peut économiser l'expansion ; pour quatre
sorties à 256 dimensions, il perd face à la récursion. Les deux sorties
impaires du lot de quatre sont structurellement nulles dans les familles
exclusivement paires. Les seuils automatiques demandent donc confirmation sur
d'autres grades, métriques et ordres de passage.
\FloatBarrier

\subsection{Jointure ternaire préparée : première comparaison}
\label{sec:tripleperf}
\paragraph{Pourquoi cet essai ?}
Une jointure de trois facteurs évite de matérialiser certains résultats
intermédiaires, mais son plan coûte à préparer. L'expérience cherche
l'horizon où cette préparation s'amortit pour les mêmes supports et sorties.
La campagne TripleJoin valide 192 cas dans les dimensions
$2,3,4,5,8,12,32,65,128,129$, pour une ou quatre sorties et des horizons
$H=1,32,1024$ ; 10\,000 appels sont ajoutés en 65 et 128D. Les
trois voies évaluées calculent les mêmes coefficients de $(ab)c$ :
deux produits complets puis extraction, \code{join3} appelé pour chaque
sortie, et \code{TripleJoinPlan} avec workspace. Les trois entrées changent
de valeurs à chaque pas mais gardent leurs supports. L'oracle \code{Int64}
indépendant vérifie les sommes de la trace ; les petits entiers du corpus
restent exactement représentables en \code{Float64}. La préparation de
l'artefact entre \emph{une fois} dans chaque trace mesurée et la sortie
est accumulée dans un tableau de l'appelant.

\begin{table}[ht]
\centering\small
\begin{tabular}{@{}rrcc@{}}
\toprule
Sorties & $H$ & Victoires direct / \code{join3} / plan & Ratio médian plan/\code{join3} \\
\midrule
1 & 1 & 2 / 8 / 0 & 14,039 \\
1 & 32 & 0 / 5 / 5 & 1,113 \\
1 & 1024 & 0 / 0 / 10 & 0,139 \\
4 & 1 & 5 / 5 / 0 & 5,427 \\
4 & 32 & 0 / 1 / 9 & 0,306 \\
4 & 1024 & 0 / 0 / 10 & 0,033 \\
\bottomrule
\end{tabular}
\caption{Première comparaison TripleJoin, médianes par cellule, préparation incluse.
Le ratio est la médiane des ratios par dimension ; inférieur à un favorise
le plan.}
\label{tab:tripleperf}
\end{table}

\begin{figure}[ht]
\centering
\begin{tikzpicture}
\begin{loglogaxis}[
  width=.82\linewidth,height=6.0cm,
  xmin=.75,xmax=1400,ymin=.02,ymax=25,
  xtick={1,32,1024},xticklabels={1,32,1\,024},
  ytick={.03,.1,1,10},yticklabels={0.03,0.1,1,10},
  xlabel={Produits réutilisant les mêmes supports},
  ylabel={Temps plan / temps \code{join3}},
  grid=major,grid style={gray!20},
  legend style={at={(.98,.98)},anchor=north east,draw=none,fill=white},
]
\addplot+[black,dashed,mark=none] coordinates {(1,1) (1024,1)};
\addplot+[blue!80!black,thick,mark=*,mark size=2.4pt] coordinates
  {(1,14.039) (32,1.113) (1024,.139)};
\addplot+[orange!85!black,thick,mark=square*,mark size=2.4pt] coordinates
  {(1,5.427) (32,.306) (1024,.033)};
\legend{Égalité,Une sortie,Quatre sorties}
\end{loglogaxis}
\end{tikzpicture}
\caption{Amortissement du plan ternaire sur le corpus TripleJoin :
médiane des ratios par dimension, préparation comprise. Sous la ligne 1,
le plan est plus rapide. Les segments guident l'œil entre les trois horizons
mesurés ; ils ne définissent aucun seuil interpolé. Temps exploratoires.}
\label{fig:tripleamort}
\end{figure}

Le plan perd au premier appel dans toutes les dimensions testées. En 65D,
une sortie et $H=32$, \code{join3} garde l'avantage (0,095 contre
0,329~ms pour la trace) ; en 8D, quatre sorties et $H=32$, le plan gagne
(0,085 contre 0,353~ms). À $H=1024$, il gagne les vingt cellules
de ce corpus, avec moins d'octets alloués \emph{cumulés} :
pour une sortie, les médianes sont environ 57,6, 5,36 et 0,178~Mo
pour direct, \code{join3} et plan ; ce ne sont pas des pics RSS.
La trace ne conserve que les sommes, pas chaque multivecteur emprunté.

Les opérandes comptent au plus quatre masques chacun et seulement deux à
quatre directions actives : 129D ne signifie pas une arithmétique dense
en 129 dimensions. PerfChecker retourne sept échantillons à $H=1$,
un à sept à $H=1024$, et parfois un seul à $H=10000$ ; les temps
froids ne disposent que d'un passage par cas. Les classements proches à
$H=32$ et les ratios de longue trace attendent donc plus de répétitions,
des supports plus riches, des métriques variées et un CPU isolé. Les
coûts froids et les refus de budget ne sont pas fusionnés avec ces temps
chauds ; aucun seuil automatique n'est encore adopté.
\FloatBarrier

\subsection{Plans, tampons et propriété des résultats}
\paragraph{Pourquoi cette technique ?}
Le \code{ProductPlan} garde les chemins. Le \code{ProductWorkspace} conserve
les tableaux de valeurs et de sortie. Ce découpage
rejoint les workspaces modulaires du calcul creux~\cite{sparseworkspace}.
Le plan convient à des supports réutilisés ; le workspace vise des appels
nombreux à valeurs changeantes, lorsque le résultat peut être consommé avant
l'appel suivant. Sa sortie est \emph{empruntée} et sera écrasée.
\paragraph{Essais, résultats et limites.}
Le remplissage, la remise à zéro, la validation des supports et la préparation
doivent être comptés ; conserver chaque sortie exige une copie. Sur le corpus
2--12D, 150 cas chauds n'allouent rien dans le noyau mesuré. À 2D creux
et $H=1$, le workspace chaud prend 0,224~\(\mu\)s mais son coût total
48,194~\(\mu\)s : la préparation renverse la lecture du seul noyau
(tableau~\ref{tab:workspace}). À 12D creux et $H=1024$, le total de
1\,506~\(\mu\)s bat les 49\,983~\(\mu\)s directs pour cette trace.
Le budget de 64~Mio à la construction ne borne ni RSS du processus, ni JIT,
ni allocations de nombres à précision arbitraire. Les prochaines comparaisons
doivent séparer résultat consommé immédiatement et historique possédé par
l'appelant, varier supports et types, et mesurer pic RSS ainsi que coût total.

\subsection{Petites dimensions et mémoire réutilisée}
\paragraph{Pourquoi cet essai ?}
Dans une petite algèbre répétée, les allocations de sortie et de
préparation peuvent dominer le calcul. L'essai mesure si réutiliser un
workspace supprime ces allocations chaudes et à quel horizon son
coût de construction est remboursé.
Une campagne distincte a exécuté 687 cas valides, dont 600 mesures d'exécution,
sur les dimensions entières 2 à 12 selon la famille et les horizons
$H\in\{1,32,1024\}$. Sur les 75 cas chauds \code{run\_product!} et les 75 cas
\code{run\_product\_values!}, l'allocation enregistrée par appel de noyau est
nulle. La préparation, l'oracle, les conversions et la conservation éventuelle
de chaque sortie changent le coût total ; la campagne a justement séparé
« chaud » et « total ».

\begin{table}[ht]
\centering\small
\begin{tabular}{@{}lrrrr@{}}
\toprule
Famille, $n$, $H$ & Direct chaud & Plan chaud & Workspace chaud & Workspace total \\
\midrule
Creux, 2, 1 & 0,496 & 2,252 & 0,224 & 48,194 \\
Creux, 2, 32 & 16,025 & 27,170 & 3,759 & 49,688 \\
Creux, 12, 1024 & 49\,983 & 8\,683 & 1\,440 & 1\,506 \\
Grade 2 plein, 12, 32 & 41\,256 & 1\,849 & 378 & 1\,251 \\
Dense, 4, 1024 & 51\,312 & 1\,724 & 527 & 585 \\
\bottomrule
\end{tabular}
\caption{Temps en microsecondes pour la trace entière. Les sorties du
workspace sont consommées sans copie de conservation.}
\label{tab:workspace}
\end{table}

Sur cette grille, le premier horizon testé où le coût total du workspace gagne
est 1024 pour la famille creuse en 2--3D, et 32 en 4--12D ; ce n'est pas une
estimation précise du seuil. Le budget de plan est de 65\,536 chemins et le
workspace de 64 Mio. Le pic RSS de la campagne atteint environ 1,54 Go,
comprenant processus et compilateur. Les gains larges et les allocations
chaudes nulles sont des constats propres au corpus, pas une garantie mémoire
générale. Les méthodes retournant directement des valeurs de coefficients
ont un contrat de sortie différent et ne sont pas assimilées aux multivecteurs.

\begin{figure}[ht]
\centering
\begin{tikzpicture}
\begin{semilogyaxis}[
 width=.82\linewidth,height=5.7cm,
 ybar,bar width=11pt,enlarge x limits=.2,
 ymin=.1,ymax=60000,
 xtick={1,2,3,4},
 xticklabels={Creux 2D \(H=1\),Creux 2D \(H=32\),Creux 4D \(H=32\),Creux 12D \(H=1024\)},
 x tick label style={rotate=18,anchor=east},
 ylabel={Temps de la trace (\(\mu\)s)},
 grid=major,grid style={gray!20},
 legend style={at={(.03,.97)},anchor=north west,draw=none,fill=white},
]
\addplot+[fill=gray!50,draw=black,mark=none] coordinates
 {(1,.520) (2,14.050) (3,141.928) (4,48053.051)};
\addplot+[fill=blue!60,draw=blue!80!black,mark=none] coordinates
 {(1,48.194) (2,49.688) (3,64.643) (4,1506.102)};
\legend{Direct total,Workspace total}
\end{semilogyaxis}
\end{tikzpicture}
\caption{Coût total de deux stratégies avec le même contrat de sortie :
le workspace inclut sa préparation une fois par trace. En 2D le coût
initial domine ; dès 4D et 32 appels sur cette famille, il est amorti.
Les sorties sont consommées sans copie. Les quatre charges ne
constituent pas une loi de croissance.}
\label{fig:workspaceamort}
\end{figure}

\paragraph{Pourquoi profiler ce workspace ?}
Les temps de la campagne précédente indiquent \emph{quand} la mémoire
réutilisée est avantageuse, mais ne localisent pas les fonctions exécutées
pendant un appel. Un premier profil PerfChecker examine donc un produit
workspace en 7D à l'horizon $H=32$, pour orienter une inspection du code,
sans transformer des piles échantillonnées en attribution causale de temps.
La première tentative, avec une opération par fenêtre de profilage, n'avait
produit aucune pile CPU. La reprise regroupe 100 épisodes de 256 opérations,
avec vérification de chaque sortie. Ses sept observations de benchmark chaud
forment une collecte \emph{distincte} de ces épisodes longs.

Dans cette reprise, les six capacités natives exécutées sont complètes :
benchmark, profils CPU et d'allocations, diagnostics de latence, de GC et de
mémoire. Le typage a été capturé, sans verdict automatique de stabilité ;
les piles CPU et d'allocations sont exportées. Sur 134 échantillons CPU,
34 piles contiennent la validation de l'algèbre et 16 celle des supports.
Ces comptages sont \emph{inclusifs} : une même pile peut contenir plusieurs
fonctions. Sous l'interférence Étendue3D, ils indiquent des endroits à
examiner, sans conclusion causale ni seuil de sélection. Le diagnostic
\emph{heap} figurait dans une autre tentative interrompue par dépassement
du budget de disque temporaire ; il n'est pas un résultat de ces six
capacités.
\FloatBarrier

\subsection{Largeur des masques et frontière 64/128}
\paragraph{Pourquoi cette technique ?}
Le masque de lame rend XOR, intersection et parité des permutations
directement calculables par opérations binaires
(équation~\eqref{eq:blade}). Garamon.jl utilise \code{UInt64} jusqu'à 64
directions, \code{UInt128} jusqu'à 128, puis \code{BigInt} ; ces types
sont fournis par Julia~\cite{julianumbers}. Un entier de
précision arbitraire reste nécessaire au-delà de 128. Le changement vise les
supports creux de dimension ambiante 65--128, où un tableau dense est
irréaliste mais le nombre de masques actifs peut rester petit.
À 65 dimensions, un masque \code{UInt64} ne suffit plus ; passer
immédiatement à \code{BigInt} peut faire payer une arithmétique et des
allocations disproportionnées pour des masques encore contenus dans
128 bits. La comparaison teste ce coût de représentation à la frontière.
\paragraph{Essais, résultats et limites.}
Depuis le premier balayage, les masques utilisent \code{UInt128} entre 65
et 128 dimensions. Les tests de correction couvrent 63, 64, 65, 70, 127,
128 et 129 dimensions. Un second PerfChecker a validé 128 cas aux
frontières 63--66 et 127--130, puis une comparaison appariée dans le même
processus a validé 24 cas de noyau diagonal simplifié. À 65D, la médiane du
coefficient ciblé est passée de 2,129~\(\mu\)s avec \code{BigInt} à
0,179~\(\mu\)s avec \code{UInt128} ; le produit entier correspondant passe
de 86,720 à 4,209~\(\mu\)s et de 159\,784 à 8\,816 octets alloués.
À 128D, les couples sont 3,179/0,204~\(\mu\)s pour le coefficient et
99,299/5,671~\(\mu\)s pour le produit entier. Ces chiffres appariés
excluent la création des supports et la première compilation ; ils ne sont
pas un gain de même amplitude démontré sur toutes les API. La voie
\code{BigInt} demeure nécessaire à partir de 129 dimensions. Il reste à
mesurer les plans, jointures, métriques et expressions de part et d'autre
de 64, 128 et 129 dimensions, en enregistrant le type \emph{effectif} du
masque, les allocations et les conversions. Les anciennes campagnes en
66/128D avec \code{BigInt} ne sont pas des répétitions d'un code identique.

\begin{figure}[ht]
\centering
\begin{tikzpicture}
\begin{semilogyaxis}[
 width=.84\linewidth,height=5.7cm,
 ybar,bar width=11pt,enlarge x limits=.25,
 ymin=.1,ymax=200,
 xtick={1,2,3,4},
 xticklabels={65D coefficient,65D entier,128D coefficient,128D entier},
 x tick label style={rotate=18,anchor=east},
 ylabel={Médiane du noyau (\(\mu\)s)},
 grid=major,grid style={gray!20},
 legend style={at={(.5,1.03)},anchor=south,legend columns=2,column sep=12pt,draw=none,fill=white},
]
\addplot+[fill=gray!50,draw=black,mark=none] coordinates
 {(1,2.129) (2,86.720) (3,3.179) (4,99.299)};
\addplot+[fill=blue!60,draw=blue!80!black,mark=none] coordinates
 {(1,.179) (2,4.209) (3,.204) (4,5.671)};
\legend{\code{BigInt},\code{UInt128}}
\end{semilogyaxis}
\end{tikzpicture}
\caption{Deux sorties du \emph{même noyau diagonal ciblé}, comparées à
chaque dimension dans un processus apparié. L'échelle est logarithmique ;
préparation des supports et compilation sont exclues. Ces résultats
ne sont pas un ratio de bout en bout pour toute l'API.}
\label{fig:masques}
\end{figure}
\FloatBarrier

\subsection{Code généré, JIT et catalogue natif}
\paragraph{Pourquoi cette technique ?}
Comme GAALOP~\cite{gaalop} pour des expressions d'algèbre géométrique, une
spécialisation peut remplacer des boucles et indices dynamiques par un
programme adapté à une structure connue. Ici, le générateur Julia déroule
au plus 256 chemins ; un plan préparé reste l'alternative quand les supports
ou types sont nombreux. La compilation Julia a son cycle de vie propre
\cite{juliaprecompile}, distinct de la construction du plan et de la
conservation de ses données.
\paragraph{Essais, résultats et limites.}
L'étude de douze familles valide 288 sorties préparées et 288 générées
contre le produit direct, mais une topologie nouvelle de 64 chemins peut
coûter environ 100~ms de JIT, et 256 chemins environ 440~ms. Les 36
décisions de coût modélisées sur ce corpus retiennent toutes le plan préparé ;
elles ne sont pas 36 traces exécutées indépendantes. Dans l'étude du
cycle de vie du catalogue, le
catalogue natif EGA3 supprime le JIT du premier produit connu après
redémarrage (tableau~\ref{tab:precompile}), mais onze topologies 4D nouvelles
ajoutent 648--658~ms de JIT à \emph{chaque} reprise. Un seul build mesuré
attribue au catalogue 5,47~Mio et 1,67~s supplémentaires. Le protocole
futur doit stratifier structures connues et inédites, mesurer directement
des épisodes de $H$ appels en processus neufs puis chauds, et inclure taille
du cache, mémoire de code et fréquence réelle des nouvelles topologies.

\subsection{Précompilation et redémarrage}
\paragraph{Pourquoi cet essai ?}
Un bon temps chaud ne garantit pas une première réponse rapide après
redémarrage. La campagne sépare le coût du catalogue, la reconstruction
des objets et la compilation du premier produit connu ou inédit.
Cette campagne compare un cache IR ou natif, chacun avec ou sans catalogue
de précompilation EGA3. Elle valide 72 processus neufs et 3\,816 observations
de phases : trois redémarrages pour chaque combinaison de quatre configurations,
deux scénarios et trois stratégies. Les algèbres, plans et tampons sont
reconstruits dans chaque processus ; un oracle indépendant vérifie les sorties.
Le tableau~\ref{tab:precompile} mesure le premier produit \emph{après}
découverte et préparation, avec sa compilation éventuelle. Il ne mesure pas
le temps depuis le lancement du processus.

\begin{table}[ht]
\centering\small
\begin{tabular}{@{}lrrr@{}}
\toprule
Configuration & Préparé & Généré & Workspace \\
\midrule
IR, sans catalogue & 64,875 & 190,480 & 62,885 \\
Natif, sans catalogue & 65,112 & 191,647 & 63,643 \\
IR, avec catalogue & 59,755 & 143,224 & 59,069 \\
Natif, avec catalogue & 0,008 & 0,122 & 0,015 \\
\bottomrule
\end{tabular}
\caption{Premier produit EGA3 après reconstruction des objets : médiane de
trois processus neufs, en millisecondes.}
\label{tab:precompile}
\end{table}

Avec catalogue natif, le compteur JIT du premier produit connu est nul dans
les neuf processus concernés. La première \emph{construction} du workspace
conserve pourtant 2,63--2,77~ms de JIT : préparation structurelle et produit
ne doivent pas être fusionnés dans une seule latence. Le lancement jusqu'au
premier résultat connu, harnais compris, passe d'environ 2,7--2,9~s sans
catalogue natif à 1,17--1,19~s avec lui. Une seule construction de cache a été
observée par configuration : le catalogue natif ajoute 5,47~Mio et environ
1,67~s à cette construction, sans établir de seuil d'amortissement.

Les douze contextes 4D découverts pendant chaque session possèdent douze
topologies distinctes, mais un seul type de plan. Après le premier contexte,
les onze suivants ajoutent 648--658~ms de JIT pour la voie générée à
\emph{chaque} redémarrage et aucun pour les voies préparée et workspace.
Les empreintes des caches disque restent inchangées dans les 72 sessions :
le JIT des nouvelles topologies n'y est pas conservé implicitement. Le
protocole ne teste ni image système, ni sérialisation des plans, ni catalogue
exhaustif ; les phases, tailles et exclusions appartiennent au protocole
de ces 72 processus neufs.
\FloatBarrier

\subsection{Longues traces et génération}
\paragraph{Pourquoi cet essai ?}
Un produit isolé ne représente pas une session qui répète des expressions
avec supports stables ou changeants. Cette campagne mesure le coût total
sur plusieurs horizons, y compris la préparation et la génération,
pour voir quelles stratégies s'amortissent réellement.
Une autre campagne exécute réellement 1, 32, 1024 ou 10\,000 itérations de
$(ab)c$ dans des supports creux de 66 et 128 dimensions, avec supports stables,
alternants ou par phases. Elle précède le passage aux masques \code{UInt128}
de 65 à 128 dimensions et décrit donc une autre révision du code. Les
coefficients changent à chaque itération, les plans sont préparés à la
première apparition d'un contexte et tous les
résultats scalaires sont vérifiés par un oracle entier. Pour 10\,000 sorties à
66D stable, les temps de trace réchauffée, préparation paresseuse comprise,
sont 1\,477 ms directs, 66,9 ms par plans, 89,7 ms générés, 101,1 ms par
\code{join3} et 27,9 ms par workspaces. À $H=1$, \code{join3} prend 0,047--0,055
ms suivant la trace, contre 1,13--1,68 ms pour la création d'un workspace.
Les traces à 10\,000 appels n'ont qu'\emph{un} échantillon retenu et la
campagne a partiellement recouvert une autre utilisation du CPU : ce sont des
résultats de criblage, sans estimation fiable de la dispersion. Les coûts
froids et de compilation ne sont pas inclus dans ces temps de trace chaude.
Ces traces répètent des produits à coefficients modifiés ; la sortie d'un
appel n'est pas l'entrée du suivant. Elles ne testent donc pas la propagation
d'erreur d'une récurrence temporelle.

Une étude séparée de douze familles de produits préparés et déroulés observe
qu'une topologie nouvelle de 64 chemins peut coûter environ 100 ms de
compilation, et de 256 chemins environ 440 ms dans les familles citées. Les
288 résultats préparés et 288 résultats générés ont égalé le produit direct
sur les cas vérifiés. Pour ces API et ce corpus, les 36 décisions prises sur
le sous-ensemble d'apprentissage pour trois horizons de réutilisation ont toutes
retenu le plan préparé ; leur coût total est modélisé, pas mesuré par une trace
distincte. Certains supports distincts partagent déjà une topologie de code.
La famille ambiante 66D de cet essai utilisait encore \code{BigInt}. Une
compression supplémentaire des programmes doit être motivée par un corpus
plus varié, un partage réel de JIT et une mesure de la mémoire de code ;
compter les seules topologies ne suffit pas.
\FloatBarrier

\subsection{Cache structurel et roulette d'éviction}
\paragraph{Pourquoi cette technique ?}
Le cache Julia conserve des \code{ProductPlan} par clé de métrique, base,
supports, opération et budget, avec LRU borné. Une roulette expérimentale
choisit une victime avec poids d'éviction $1/(1+h)$ pour $h$ succès ;
elle ne tire \emph{aucune} contribution algébrique. TinyLFU
\cite{tinylfu} et SIEVE~\cite{sieve} proposent d'autres politiques de cache
dans leurs domaines d'origine, pas des accélérations prouvées pour Garamon.
Le cache vise des structures récurrentes quand l'appelant ne garde pas
directement chaque plan.
\paragraph{Essais, résultats et limites.}
Sur 40 produits avec deux supports chauds et un intrus, les temps direct,
LRU et roulette sont 1,60, 2,42 et 1,83~ms ; sur un balayage ils prennent
1,56/3,11/3,15~ms. Une trace à quatre phases favorise les deux caches
(1,52/0,736/0,763~ms), tandis qu'un remplacement brutal du jeu chaud
favorise LRU (4,09/1,11/4,08~ms), avec 76 succès sur 80 contre 43
pour la graine de roulette rapportée. La politique par
défaut reste donc LRU et l'accès au cache explicite. Pour décider d'une
admission nouvelle, il faut mesurer de vraies traces de supports, coût des
clés et verrous, octets retenus, taux de succès, coût de régénération et
variation entre graines. Un cache de \emph{valeurs} aurait un autre contrat
d'invalidation et n'est pas implémenté.

\subsection{Formes factorisées, sous-espaces actifs et trains tensoriels}
\paragraph{Pourquoi cette technique ?}
Une lame simple peut rester un produit extérieur de vecteurs ; des réflexions
peuvent rester une chaîne de facteurs ; une sous-algèbre de directions
coordonnées actives réduit la dimension effectivement parcourue. Les trains
tensoriels de Oseledets~\cite{oseledets} représentent un tenseur de
coefficients par cœurs successifs. Pour le produit de Clifford \emph{diagonal},
la construction Julia conserve une parité des bits droits déjà lus, multiplie
le facteur $(-1)^{a_i p}g_i^{a_i b_i}$ et forme un train exact de rang
non comprimé au plus $2r_A r_B$. La borne de rang est une dérivation
spécifique à cette construction, pas un théorème attribué à
Oseledets. Les seuils de rang et d'entrées refusent une expansion hors budget.
Ces voies visent des entrées déjà structurées, pas des coefficients arbitraires.
\paragraph{Essais, résultats et limites.}
Les 84 vérifications TT incluent rangs 1 et 2 et 66D. Avec 12 bits actifs
à 66D, chaque entrée explicite possède 4\,096 coefficients ; le TT de
résultat avec algèbre occupe 15\,757 octets contre environ 380\,021
pour une entrée explicite. Pour une sortie de grade 6, un cas mesure
42,37~\(\mu\)s de produit TT et 6,86 de lecture, contre
8\,656,07~\(\mu\)s pour la jointure de supports explicites.
Cette comparaison commence avec les entrées \emph{déjà séparables} :
la conversion d'une entrée arbitraire vers un petit rang pourrait être
impossible ou plus coûteuse. Il n'y a ni troncature approximative, ni produit
TT non orthogonal validé. Le test futur apparie la même sortie en comptant
construction des facteurs ou cœurs, aller-retour, rang atteint, mémoire,
conditionnement et éventuelles contraintes de dualité ambiante.

\subsection{ZDD et changement de représentation}
\paragraph{Pourquoi cette technique ?}
Les diagrammes de décision à zéros supprimés de Minato~\cite{minato}
compressent des \emph{familles d'ensembles}. Comme les masques actifs
forment une famille de sous-ensembles des directions, un ZDD peut décrire
un support régulier sans lister toutes ses lames. Le prototype initial
ne stockait que l'appartenance et la cardinalité. Un nouveau ZDD pondéré
associe des rationnels exacts aux terminaux ; son produit Clifford récursif
mémoïse les paires de sous-diagrammes et la parité des bits du second
opérande. Pour une métrique diagonale $g$, les branches où une même
direction est présente dans les deux opérandes portent le facteur $g_i$ ;
les branches résultantes de même masque sont additionnées exactement.
L'algorithme ne supprime aucune contribution non nulle, et impose des
budgets explicites de nœuds, de travail et de taille des coefficients.
Les représentations matricielles ou les
changements vers un sous-espace linéaire actif restent des hypothèses :
la conversion, le conditionnement et la sortie en base d'origine pourraient
annuler tout gain.
\paragraph{Essais, résultats et limites.}
Sur 10\,626 masques de grade 4 en 24D, le support régulier forme
84 nœuds et 3\,272 octets contre 15\,138 nœuds et 659\,216 octets pour
un support aléatoire de même cardinal ; leurs constructions prennent
5,44 et 17,65~ms. Ce résultat est une comparaison de
\emph{formats de support}, pas de produits. Le nouveau produit pondéré passe
un prévol rationnel sur une métrique signée et dégénérée en 4D, avec un
oracle indépendant par mots de Clifford ; 32 assertions du noyau et six
du raccord ont aussi été validées. La compression n'est pas garantie :
sur le cas de grade 2 en 6D du banc exploratoire, le résultat compte
234 nœuds pour seulement 28 coefficients non nuls. La prochaine mesure
devra comparer ZDD, dictionnaire, liste et blocs sur les mêmes valeurs,
en comptant construction, calcul et extraction. Les cas irréguliers
fournissent le témoin négatif indispensable.

\subsection{Élagage approché, résidus et replay}
\paragraph{Pourquoi cette technique ?}
Une contribution peut être omise sous seuil, ou acceptée avec une probabilité
et repondérée ; un journal de résidus peut ensuite corriger une sortie, ou
un segment être rejoué depuis un état exact. Les estimateurs à correction
aléatoire de Rhee et Glynn~\cite{rhee} motivent la question du biais,
sans garantir la précision d'une trajectoire non linéaire. Ces méthodes
exigent un \emph{mode approché explicite} et une règle d'erreur ; elles
ne sont pas des routes cachées du produit exact.
\paragraph{Essais, résultats et limites.}
Le prototype de 2\,700 cas forme encore toutes les contributions candidates :
il ne démontre aucune économie de parcours. Le seuil donne un ratio médian
de temps 0,817 en bilinéaire avec erreur, mais 1,017 en affine ; le
résidu corrige exactement les 540 sorties rationnelles finales examinées,
au prix de ratios 1,149/1,433, et ne retrouve les bits \code{Float64}
que dans 279 cas. Le replay récupère les 540 témoins flottants bit à bit,
avec ratios 1,140/1,198 (tableau~\ref{tab:pruning}). Un tirage Bernoulli
non biaisé à un pas ne borne pas les erreurs rares après récurrence.
L'étape de recherche utile est d'éviter \emph{avant formation} des blocs
entiers par bornes certifiées, d'enregistrer les résidus et de mesurer erreur,
queue, coût de certificat et correction sur des horizons plus longs,
toujours contre un témoin exact de même sortie.

\subsection{Élagage approché dans des récurrences : résultat exploratoire}
\paragraph{Pourquoi cet essai ?}
Omettre de petites contributions peut accélérer un calcul, mais leur
effet se propage dans une récurrence. L'expérience confronte le coût et
l'erreur de l'élagage à ceux d'une restauration exacte de la trajectoire.
Un prototype \emph{séparé du noyau public exact},
a validé 2\,700 cas et refusé explicitement 540 cas supplémentaires avant
exécution. Ici, la validation porte sur le protocole et les assertions de
chaque méthode ; elle ne signifie pas que les sorties approchées sont exactes.
Il parcourt douze dimensions de 2 à 128, avec au plus quatre
directions actives, trois classes de métriques, récurrences affines et
bilinéaires, et 7 échantillons temporels par cas. Les dimensions ambiantes
élevées n'augmentent donc pas ici la taille arithmétique au-delà des quatre
directions actives. Toutes les méthodes forment encore chaque contribution
candidate ; elles diffèrent dans la mise à jour ou la correction des sorties.

\begin{table}[ht]
\centering\small
\begin{tabular}{@{}lrr@{}}
\toprule
Méthode & Ratio médian affine & Ratio médian bilinéaire \\
\midrule
Seuil de contribution & 1,017 & 0,817 \\
Roulette Bernoulli ($p=1/2$) & 1,115 & 0,904 \\
Résidu transporté et corrigé & 1,433 & 1,149 \\
Replay exact du segment & 1,198 & 1,140 \\
\bottomrule
\end{tabular}
\caption{Temps total de la méthode divisé par celui du témoin exact de la même
configuration ; préparation commune exclue, correction incluse. Médiane des
ratios, pas ratio des médianes, dans le corpus de récurrences décrit ici.
Une valeur inférieure à un indique un gain de temps observé.}
\label{tab:pruning}
\end{table}

Le seuil bilinéaire réduit donc de 18,3\% la médiane des ratios de temps sur ce
seul protocole, avec une erreur. Le résidu restaure exactement la sortie finale
en rationnels dans 540 configurations, mais ne retrouve le témoin flottant
bit à bit que dans 279 ; le replay retrouve les 540 sorties flottantes bit à
bit, au coût médian supérieur. La roulette de probabilité $1/2$ est
conditionnellement non biaisée à un pas, mais peut produire des trajectoires
rares d'erreur énorme ; 32 graines par configuration ne suffisent pas à en
mesurer correctement la queue. Aucun résultat ne démontre un gain du
pruning sur \code{ProductWorkspace}, ni une voie exacte plus rapide. Les
erreurs intermédiaires et le coût de reprise appartiennent au même
contrat de comparaison que la sortie finale.
\figureattente{Roulette : balayage de la probabilité de conservation}
{Abscisse : probabilité $p$ et horizon ; ordonnées : coût total, erreur,
variance et queue d'erreur face au témoin exact.}
{Étude paramétrique de la roulette au-delà du seul point $p=1/2$ du
tableau~\ref{tab:pruning}.}{fig:roulette-parametres}
\FloatBarrier

\subsection{Sélecteur contextuel et réglage des hyperparamètres}
\paragraph{Pourquoi cette technique ?}
SATzilla~\cite{satzilla}, ASlib~\cite{aslib} et
AutoFolio~\cite{autofolio} formalisent la sélection d'algorithme et
la configuration d'un sélecteur. Pour Garamon, les variables disponibles
\emph{avant} le choix comprennent dimension, type de masque, métrique,
tailles et grades des supports, sorties demandées, horizon annoncé et
présence réelle d'un plan ou workspace compatible. Le choix reste parmi
des stratégies exactes \emph{admissibles} ; extraire des caractéristiques,
inférer, vérifier et se replier ont un coût. SHAP~\cite{shap} peut décrire
les contributions d'un modèle entraîné, sans valider sa causalité ni son
gain.
Le futur candidat de sélection sera un assemblage de stratégies
mathématiquement compatibles \emph{avec} son vecteur de paramètres :
budget de cache, seuil récursif, taille de workspace ou de lot, rang de
représentation, nombre de threads et politique de sortie selon le cas.
Les valeurs par défaut, limites et interactions seront d'abord contrôlées
par oracle puis explorées dans des grilles DrWatson ; l'optimisation ne
commencera qu'après les mesures isolées. Le contrat exact et les budgets
RAM/VRAM sont des contraintes d'admission, non des pénalités compensables
par quelques microsecondes gagnées. La roulette approximative forme un
espace distinct : sa probabilité de conservation sera balayée, avec
erreur, variance et temps à chaque valeur, sans mélange avec les routes
exactes. La sélection finale sera jugée sur des familles et des supports
tenus à l'écart, coût de décision compris.
\paragraph{Premier prototype et mesure.}
Une règle analytique \emph{non ajustée} choisit maintenant parmi les cinq
routes exactes pour les coefficients demandés de $(ab)c$ ; ses poids sont
des hypothèses, pas des durées apprises. Les
406/406 tests ciblés passent dans l'état actuel. Un premier smoke
PerfChecker à un thread, en
EGA4 euclidienne, un coefficient scalaire et $H=1$, valide les six routes
(cinq forcées et le choix automatique), avec 15 observations chacune,
soit 90 échantillons. L'épisode chaud comprend extraction, décision,
préparation et sortie possédée. La règle choisit \code{join3} :
60,135~\(\mu\)s en médiane contre 59,466~\(\mu\)s pour la même route forcée.
Ces séries ont tourné sous activité Étendue3D ; l'écart de
0,669~\(\mu\)s ne peut être attribué au seul routeur.
Les sondes de première compilation sont diagnostiques et non additives :
elles ne sont pas incluses dans ces médianes chaudes.

La matrice préparée compte 144 structures et 5\,184 identités de replay,
mais seulement huit groupes d'apprentissage indépendants après purge,
contre un minimum déclaré de vingt. Aucun HPO, regret sur familles
tenues à l'écart, seuil de routage ni gain général n'est donc validé.
Les anciennes campagnes de produit triple ciblé et de produit binaire
complet ne constituent pas des étiquettes comparables pour l'entraînement.
Le protocole figé prévoit la même expression
$(ab)c$, une
ou quatre sorties appartenant à l'appelant, $H=1,32,1024$, et cinq
candidats exacts : complet, récursif, \code{join3}, préparé et workspace.
Chaque coût d'épisode \emph{à qualifier} distinguera admission,
préparation, premier JIT, exécutions,
reconstructions dues aux annulations et copies. Les familles d'algèbres
\emph{et} les familles de supports sont disjointes entre apprentissage,
validation et test ; les renommages d'un support restent groupés.
Les témoins sont direct, meilleur choix constant et règle analytique ;
un petit arbre ou modèle de coûts régularisé reçoit au plus douze
configurations choisies sur validation. Le résultat à publier est le
coût total, la mémoire et le regret face au meilleur candidat
rétrospectif mesuré, avec refus et queue des pertes. L'oracle rétrospectif
n'est pas une politique exécutable gratuitement.
\figureattente{Sélection, HPO et explications}
{Panneaux prévus : regret sur familles tenues à l'écart, coût du sélecteur,
sensibilité aux paramètres et attribution SHAP du modèle retenu.}
{Évaluation après acquisition de groupes d'apprentissage indépendants
et campagne isolée.}{fig:hpo-shap}

\subsection{Threads, processus, lots compacts et GPU}
\paragraph{Pourquoi cette technique ?}
Des produits indépendants peuvent partager un plan immuable et attribuer
un workspace mutable à chaque tâche. Le manuel Julia~\cite{juliathreads}
impose de prévenir les accès concurrents aux données mutables ; une tâche
peut migrer, donc un tampon indexé naïvement par numéro de thread ne suffit
pas. Des processus ajoutent démarrage, sérialisation et duplication des
états. Un GPU demande un véritable noyau de contraction de chemins,
avec données compactes et sorties définies ; les recommandations
CUDA.jl~\cite{cudaperf} rappellent de mesurer transferts et lancement.
\paragraph{Essai actuel et expérience à conduire.}
Le format \code{PackedProductBatch} implémenté donne, pour 32 répétitions des
\emph{mêmes objets} en 12D, 39,93~\(\mu\)s tout compris contre
71,62~\(\mu\)s pour le lot préparé ordinaire ; avec 32 objets distincts
de mêmes supports, il perd (586,06 contre 72,62~\(\mu\)s)
. Cela isole déjà le risque de conversion ; aucun gain
parallèle ou GPU n'en découle. Un protocole CPU préparé
 compare 1/2/4/8/16 threads et 1/2/4/8 processus
sur lots $B=1,32,1024$, même oracle et sortie possédée par l'appelant.
Il distingue démarrage, préparation, premier JIT, exécution, copie,
sérialisation et RSS ; aucune accélération n'est encore validée.
Les points à huit et seize threads du protocole complet attendent une
machine isolée ; l'exploration sous charge concurrente reste bornée à
quatre threads ou workers.
Pour le GPU local de 8~Gio de VRAM, le futur comparatif doit ajouter
empaquetage, transferts aller-retour, noyau \emph{synchronisé}, mémoire active,
précision et cas à données déjà résidentes. La taille du lot doit amortir
ces coûts avant toute revendication.

\paragraph{Premier prévol CUDA, sans classement de vitesse.}
Un adaptateur optionnel exécute maintenant le produit packed exact sur CPU et
sur CUDA à partir du même plan et restitue, dans les deux cas, une matrice
possédée sur l'hôte. L'oracle entier à mots de base vérifie chaque coefficient,
indépendamment des deux noyaux. Huit cas bornés combinent les dimensions 8 et
65, les horizons 1 et 32, une signature mixte et les deux voies. Les huit
cas ont passé l'oracle et PerfChecker ; un rejeu a vérifié les archives et
ignoré les cas achevés. Un essai sans GPU a conservé la voie CPU tout en
refusant explicitement la voie CUDA. Pour chacune des quatre identités GPU,
trois traces CUPTI distinguent noyau résident, calcul résident avec retour
hôte et épisode complet avec transferts. Les instantanés de mémoire libre
encadrent ces traces, mais ne mesurent pas le pic de VRAM. Le profilage
instrumenté et les trois échantillons chauds par cas ne fondent aucun
classement de vitesse ; dimensions, signatures et horizons doivent encore
être étendus sur une machine dédiée.
\figureattente{Parallélisme CPU et calcul GPU}
{Abscisse : taille du lot et nombre de threads, processus ou blocs GPU ;
ordonnées : temps complet, temps noyau, transferts, RSS et VRAM.}
{Comparaison isolée des contrats CPU et GPU, avec résultats possédés
identiques et transferts explicitement comptés.}{fig:parallele-gpu-isole}

\section{Techniques structurelles : pourquoi les tester ?}
\label{sec:transpositions}
Les quinze mécanismes suivants répondent chacun à une situation de calcul
précise. Pour les prototypes mesurés, le principe précède immédiatement le
protocole, la figure et l'interprétation. Pour les autres, l'essai décrit une
\emph{question ouverte} : aucun chiffre de performance Garamon ne lui est
attribué. Le tableau suivant oppose les trois niveaux de travail que le
lecteur pourrait confondre : masques, coefficients et instructions. Il
situe aussi les autres prototypes mesurés, sans les classer entre eux.

\begin{table}[ht]
\centering\small
\begin{tabularx}{\linewidth}{@{}p{30mm}X X@{}}
\toprule
Technique et état & Situation qui motive l'essai & Comparaison pertinente \\
\midrule
N01, rang binaire & Beaucoup de dimensions ambiantes, mais peu de directions de masque indépendantes. & Coût total du support compact contre parcours creux et plan préparé, selon la réutilisation. \\
N03, radical & Une métrique dégénérée peut ramener l'inverse à un quotient plus petit. & Inverse par quotient et série finie contre résolution régulière complète. \\
N04, identité courte & La même entrée admet $A^2=\lambda1$ et sert à de nombreuses puissances. & Certifier puis réemployer contre puissances directes ou binaires, repli compris. \\
N05, Pfaffien & Seul le scalaire d'une chaîne de vecteurs est demandé. & Contractions et Pfaffien contre témoins de même sortie admissibles. \\
Superoptimisation, à tester & Les mêmes chemins passent par de courtes boucles de signe et d'indexation très répétées. & Instructions équivalentes sur le même noyau, puis temps total et taille du code. \\
Multimodulaire, à tester & Les grands coefficients exacts rendent l'oracle ou le calcul entier coûteux. & Modules et reconstruction contre \code{BigInt} ou rationnels directs, preuve d'exactitude comprise. \\
\bottomrule
\end{tabularx}
\caption{Pourquoi ces techniques ont été retenues et quel témoin permettrait
de juger leur intérêt. Les quatre premières ont des essais Garamon décrits
ci-dessous ; les deux dernières n'ont encore aucune mesure Garamon.}
\label{tab:questionsstructurelles}
\end{table}

Le rang binaire, le radical nilpotent et le Pfaffien disposent de
prototypes hors du cœur Julia et de mesures exploratoires
(sections~\ref{sec:n01perf}, \ref{sec:n03perf} et~\ref{sec:n05perf}) ;
leur motivation, leur protocole et leurs observations sont réunis
ci-dessous. N04, fondé sur une identité polynomiale courte, possède aussi
un prototype, des tests de correction et un premier smoke de vitesse
exploratoire.
Les onze autres mécanismes ne sont ni implémentés ni mesurés dans
Garamon.jl : leur texte explique l'usage visé et l'essai à conduire,
sans présenter ce dernier comme un résultat. Leurs références
établissent un résultat dans leur domaine d'origine, pas un gain transféré.
Chaque essai devra conserver le contrat de sortie : multivecteur complet,
projection, scalaire ou simple décision. Les économies algorithmiques
annoncées ci-dessous désignent le travail que l'on pourrait supprimer,
avant les coûts de détection, de conversion, de préparation et de mémoire.

\paragraph{Conditions communes des mesures.}
Les mesures ci-dessous ont été produites sur les scénarios et oracles
définis en section~\ref{sec:protocol}. Elles emploient Julia 1.13.0 et
un Intel Core i7-12700 ; la plupart utilisent
PerfChecker, tandis que les études de génération et de précompilation ont leurs
propres harnais. Chaque
famille a son protocole ; ni la répétition ni les intervalles statistiques ne
sont uniformes. Les temps de campagnes différentes ne forment pas un même
benchmark apparié. D'autres processus de calcul Étendue3D ont partagé la
machine pendant cette phase ; les temps sont donc des \emph{résultats
exploratoires} à confirmer par une répétition intégrale isolée. Les comptes
de cas validés décrivent la correction observée dans chaque banc, sans
transformer les classements de temps en seuils définitifs.
Dans les campagnes de dimensions, de workspace, de
précompilation et de longues traces, des oracles indépendants comptent les inversions des indices
de base avec de petits entiers ; les comparaisons \code{Float64} de ces
fixtures sont exactes dans leurs bornes de somme. L'étude de génération
compare ses sorties au produit direct de Garamon.jl. La campagne d'élagage
utilise un oracle rationnel et relève séparément l'erreur flottante. Ces
faits ne valent pas pour tous les produits flottants.

\subsection{Rang binaire des supports}
\paragraph{Pourquoi ce choix ?}
Une dimension ambiante élevée n'implique pas que les produits utilisent
toutes ses directions. Le \emph{rang binaire} $d$ compte les masques
indépendants nécessaires pour engendrer par XOR les masques rencontrés :
par exemple, deux masques indépendants n'engendrent que quatre masques,
même dans une algèbre à 128 directions. C'est pourquoi N01 a été choisi
pour des expressions à support restreint : réduire l'adressage et le
stockage possibles de $2^n$ à $2^d$. L'expérience doit vérifier si cette
économie théorique paie la construction du plan, ses contrôles et les
sorties effectivement demandées.
Les chaînes de Pauli sont codées par des vecteurs binaires et traitées par
XOR et opérations sur bits~\cite{paulib}. Dans une base \emph{orthogonale},
la relation~\eqref{eq:blade} donne également un masque de sortie $A\oplus B$.
Si les masques présents dans une expression engendrent un sous-espace de rang
$d$ sur $\mathbb F_2$, ses produits ne peuvent atteindre que $2^d$ masques,
contre $2^n$ dans l'algèbre ambiante. Le signe de permutation et les facteurs
de métrique demeurent nécessaires : cette fermeture n'est pas une
identification automatique avec $\Cl(d,0)$.


Le premier prototype N01 compare des produits de deux facteurs avec
supports fixés et sortie complète. Le cas encore visé est une longue
expression dont les supports se renouvellent dans un petit espace binaire.
Il faudra propager un certificat de base tant que les opérations préservent
cet espace, puis comparer représentation compacte et courante à support,
métrique et sortie identiques. Un cas négatif fera croître le rang à chaque
étape. Une métrique non orthogonale nécessite d'abord une preuve de
fermeture propre à cette base.

\paragraph{Essais N01 : épisodes complets observés.}
\label{sec:n01perf}
La question mesurée est simple : sur combien de produits réutilisant les
mêmes supports N01 coûte-t-il moins qu'un plan déjà préparé ou qu'un
parcours creux ? La figure~\ref{fig:n01horizon} résume cette comparaison
\emph{sur un épisode entier}, construction et conservation des sorties
comprises. La borne $2^d$ motive le test, mais ne prédit pas son temps.
Le prototype N01 reste expérimental, hors du cœur et du
sélecteur. Il réduit l'\emph{adressage} du support dans
une base orthogonale en conservant le masque, le grade, le signe et le
facteur métrique ambiants ; il parcourt toujours les paires d'entrée.
Ses 676 assertions de correction et de contrat passent. Une campagne
PerfChecker à un thread valide 648/648 cas : deux passages avec ordre
des méthodes inversé, 12 dimensions de 4 à 128, trois familles de
supports, trois horizons et trois stratégies (creux actuel, plan
préparé, N01). Les métriques diagonales positives, indéfinies et
dégénérées alternent selon la dimension ; elles ne forment donc pas un
croisement exhaustif dimension--signature. L'oracle \code{Int64}
indépendant vérifie tous les coefficients pour de petits entiers
exactement représentables en \code{Float64}.

Chaque \emph{épisode chronométré} construit un nouvel artefact, exécute
$H$ produits et conserve les $H$ multivecteurs de sortie possédés. Le
coût total par produit est la médiane des épisodes \emph{observés}
divisée par $H$. Il ne résulte pas d'une addition de médianes de
construction et d'exécution. Les 9\,396 échantillons bruts archivent
séparément 1\,512 constructions, 3\,942 exécutions chaudes et 3\,942
épisodes complets. L'ancien criblage additif de 324 cas demeure
distinct et ne doit pas être interprété comme une mesure bout en bout.
La première compilation est documentée séparément ; elle ne constitue
pas ici un démarrage de processus isolé.

\begin{figure}[ht]
\centering
\begin{tikzpicture}
\begin{axis}[
 width=.77\linewidth,height=5.6cm,
 ybar,bar width=12pt,enlarge x limits=.35,
 ymin=0,ymax=37,ytick={0,6,12,18,24,30,36},
 xtick={1,2,3},xticklabels={1,32,1\,024},
 xlabel={Produits par épisode \(H\)},
 ylabel={Cas favorables à N01 (sur 36)},
 grid=major,grid style={gray!20},
 legend style={at={(.5,1.03)},anchor=south,legend columns=2,column sep=12pt,draw=none,fill=white},
]
\addplot+[fill=blue!60,draw=blue!80!black,mark=none] coordinates
 {(1,35) (2,24) (3,0)};
\addplot+[fill=orange!65,draw=orange!85!black,mark=none] coordinates
 {(1,35) (2,23) (3,0)};
\legend{Passage 1,Ordre inverse}
\end{axis}
\end{tikzpicture}
\caption{N01 comparé au \emph{même plan préparé} sur le coût total
directement observé, construction et sorties conservées comprises.
À $H=32$, 23 configurations gagnent dans les deux passages. Les
barres nulles à $H=1\,024$ signifient 0/36 dans chaque passage,
pas des observations manquantes. Criblage exploratoire sous Étendue3D,
sans seuil de routage déduit.}
\label{fig:n01horizon}
\end{figure}
\FloatBarrier

N01 a une médiane totale inférieure à celle du plan préparé dans
35/36 configurations à $H=1$ dans chacun des deux passages, 24/36 puis
23/36 à $H=32$ (23 communes), et 0/36 à $H=1\,024$. Face au chemin
creux actuel, les comptes sont 31/36 à $H=1$ et 36/36 aux deux
horizons plus longs dans chacun des passages. Le tableau de sortie de
$2^d$ masques et les facteurs pour toutes les paires d'entrée peuvent
expliquer le recul face au plan à grand horizon ou rang plus élevé ;
ce diagnostic reste une hypothèse de coût, non une causalité isolée.
L'ordre inversé ne neutralise pas la concurrence Étendue3D : ces
classements demandent un rejeu stable avec échantillons et conditions
identiques. Le manifeste étendu de 10\,287 cas, couvrant chaque
dimension entière de 2 à 128, est \emph{préparé mais non exécuté}.

\paragraph{Premier lot du rejeu K1v2.}
Une campagne distincte confronte le plan exact préparé du cœur, N01,
un plan de rang compact et sa variante avec workspace, à entrées et
sorties possédées communes. Son manifeste K1v2 prévoit
41\,148 cas : 127 dimensions de 2 à 128, trois familles, trois métriques,
trois horizons, trois requêtes de sortie et quatre routes. Il est
partitionné en 127 lots disjoints de 324 cas. Le premier lot
\code{n002}, limité à la dimension 2 et à un seul passage, a qualifié
324/324 cas en 168,08\,s, avec 5\,292 échantillons bruts et un pic RSS
de 0,768\,Gio. Les sorties ont été vérifiées contre l'oracle avant et
après mesure. L'audit global établit 324 cas démarrés, mesurés et
qualifiés sur 41\,148, sans doublon ; les 40\,824 autres, des lots
\code{n003} à \code{n128}, ne sont pas exécutés.
La condition enregistrée est \code{exploratory\_interference} avec le
label \code{Etendue3D}. Ce lot vérifie la partition et le rejeu,
mais ne détermine ni classement stable des quatre routes ni seuil de
sélection. Ses temps ne sont pas comparés à ceux de la campagne N01
précédente.

\FloatBarrier
\subsection{Secteurs invariants}
\paragraph{Pourquoi ce choix ?}
Le \emph{tapering} des Hamiltoniens de Pauli exploite des symétries
commutantes pour restreindre le calcul à un secteur~\cite{bravyitaper}.
Pour des involutions commutantes $S_j$, on peut former
$P_{\boldsymbol\epsilon}=\prod_j(1+\epsilon_j S_j)/2$.
Une projection $P_{\boldsymbol\epsilon}AP_{\boldsymbol\epsilon}$ peut alors
être calculée dans un bloc plus petit \emph{si} tous les opérateurs de la
trace préservent ce bloc. Le gain possible porte sur les secteurs écartés ;
une sortie multivectorielle complète ne permet pas de les effacer.

\paragraph{Expérience à conduire ; aucune mesure Garamon.}
Le prototype devra certifier commutation, involution, invariance et secteur
demandé, puis comparer la projection obtenue à celle d'un calcul complet.
Des termes qui couplent les secteurs, ou une sortie portant sur tous les
secteurs, servent de contrôles négatifs. La construction des projecteurs et
la conversion des données entrent dans le coût total.

\subsection{Filtration nilpotente des métriques dégénérées}
\paragraph{Pourquoi ce choix ?}
Une métrique dégénérée ajoute des directions radicales, mais leur
nilpotence borne la série nécessaire pour inverser un élément dont le
quotient est inversible. L'essai demande si résoudre d'abord ce
quotient plus petit réduit le coût d'une inversion complète une fois
la construction, la série et les sorties comptées.
Dans une algèbre dégénérée, le radical $R$ de la forme quadratique génère
un idéal nilpotent~\cite{ablamowicz,bamberg}. Si $\dim R=r$, l'idéal
$J$ des termes portant au moins une direction radicale vérifie
$J^{r+1}=0$ et le quotient vaut $\Cl(V)/J\simeq\Cl(W)$ pour un
complément non dégénéré $W$. Une filtration par degré radical peut
donc arrêter les produits au-delà de $r$ facteurs radicaux.
Pour $A=q+j$, $j\in J$, il faut \emph{d'abord} établir que l'image
$q$ est inversible dans ce quotient et calculer réellement $q^{-1}$.
Alors $A=q(1+q^{-1}j)$ possède l'inverse fini
\[
 A^{-1}=\sum_{k=0}^{r}(-q^{-1}j)^k q^{-1}.
\]
L'ordre des facteurs est essentiel dans l'algèbre non commutative ;
un terme radical ne sauve pas un quotient singulier.

Le domaine visé est une métrique \emph{réellement} dégénérée, notamment des
formes de PGA ; un vecteur nul de CGA ne suffit pas à prouver l'existence
d'un radical. Une matrice non diagonale à diagonale nulle peut même être
non dégénérée. Le prototype N03 accepte uniquement une métrique déjà
diagonale, vérifie l'inverse du quotient, puis contrôle les deux produits
$AA^{-1}=A^{-1}A=1$ avec un oracle rationnel
indépendant. Son smoke, son écran étendu et leurs limites
sont décrits immédiatement après ce principe. Les changements de base, la taille des
coefficients et les produits de la série peuvent effacer le gain du
quotient plus petit.

\paragraph{Essais N03 : inverse conditionnelle.}
\label{sec:n03perf}
Le prototype N03 compare l'inverse par matrice régulière gauche dans
la sous-algèbre active entière à la résolution \emph{effective} de
l'inverse du quotient non dégénéré, suivie de la série finie du
radical. Les deux voies construisent leurs entrées à chaque
itération et rendent la même suite de multivecteurs inverses possédés.
L'oracle rationnel indépendant vérifie tous les coefficients et les
produits à gauche et à droite. Les 114 assertions passent, y compris
des quotients singuliers refusés et des vecteurs isotropes qui ne sont
pas dans le radical. Soixante-douze contrôles complémentaires vérifient
la grille, sa partition et la correction de la métrique signée à une
direction non radicale.

Le smoke PerfChecker valide 6/6 routes avec 42 échantillons, sept par
route : dimensions 3, 4 et 5, respectivement un, deux et trois axes
radicaux, deux axes non radicaux, horizon un. Les métriques du pilote
sont déjà diagonales ; aucune décomposition d'une métrique générale
n'est chronométrée. Toutes les entrées mesurées sont inversibles,
avec des coefficients \code{Float64} comparés à l'oracle rationnel
selon la tolérance annoncée ; leurs quotients ont une partie scalaire
dominante. Les médianes et p95 descriptifs en
microsecondes sont :

\begin{table}[ht]
\centering\small
\begin{tabular}{@{}crrrr@{}}
\toprule
$n/r$ & Régulier méd. & Radical méd. & Régulier p95 & Radical p95 \\
\midrule
$3/1$ & 19,252 & 15,830 & 44,025 & 43,590 \\
$4/2$ & 35,469 & 26,041 & 113,569 & 59,482 \\
$5/3$ & 60,028 & 39,888 & 193,544 & 73,751 \\
\bottomrule
\end{tabular}
\caption{Premier smoke N03, un thread et ordre régulier puis radical,
avec construction et sortie incluses. Sept échantillons par route sous
interférence Étendue3D : les p95 sont dispersés et aucun gain
reproductible n'est établi.}
\label{tab:n03smoke}
\end{table}

L'écran étendu corrigé couvre douze dimensions de 2 à 128, des rangs
radicaux $r=1,2,3$, une ou deux directions non radicales actives
$s$, des quotients positifs ou signés, trois horizons et les deux
méthodes. Deux processus frais, avec ordre de passage inversé, valident
chacun 756/756 cas et 5\,292 échantillons bruts : 1\,512 cas et
10\,584 échantillons au total. Les temps des
contrôleurs sont 154,49 et 155,73\,s ; leurs pics RSS, 728,62 et
715,46\,Mio. L'audit des qualifications natives, des échantillons et
de la couverture ne trouve ni manque ni doublon. L'écart absolu maximal
avec l'oracle rationnel est $5,55\times10^{-17}$. Le premier essai,
archivé séparément, avait 396 cas réussis et 360 erreurs du générateur
pour $s=1$ ; il est entièrement exclu de ces chiffres après correction
et tests de cette branche.

\begin{figure}[ht]
\centering
\begin{tikzpicture}
\begin{axis}[
 width=.82\linewidth,height=5.6cm,
 ybar stacked,bar width=15pt,enlarge x limits=.14,
 ymin=0,ymax=80,ytick={0,20,40,60,80},
 xtick={1,2,3,4,5,6},
 xticklabels={$1/1$,$1/2$,$2/1$,$2/2$,$3/1$,$3/2$},
 xlabel={Directions actives \(r/s\) : radicales / non radicales},
 ylabel={Configurations appariées},
 grid=major,grid style={gray!20},
 legend style={at={(.5,1.03)},anchor=south,legend columns=3,column sep=8pt,draw=none,fill=white},
]
\addplot+[fill=blue!60,draw=blue!80!black,mark=none] coordinates
 {(1,0) (2,20) (3,20) (4,53) (5,59) (6,54)};
\addplot+[fill=orange!65,draw=orange!85!black,mark=none] coordinates
 {(1,70) (2,34) (3,30) (4,0) (5,0) (6,0)};
\addplot+[fill=gray!55,draw=gray!75!black,mark=none] coordinates
 {(1,2) (2,12) (3,16) (4,7) (5,1) (6,0)};
\legend{Radical,Régulier,Sensible ou égalité}
\end{axis}
\end{tikzpicture}
\caption{Écran N03 : 378 configurations appariées entre deux ordres.
La couleur désigne la méthode dont la médiane d'épisode est plus petite
dans \emph{les deux} ordres. Les 38 autres configurations changent de
choix avec l'ordre ou sont à égalité. Comparaison exploratoire à un
thread sous interférence Étendue3D ; les familles et horizons sont
agrégés à l'intérieur de chaque groupe \(r/s\).}
\label{fig:n03screen}
\end{figure}

Au total, la voie radicale est plus rapide dans les deux ordres pour
206/378 configurations, la résolution régulière pour 134/378 ; les
38 autres sont sensibles à l'ordre ou à une égalité. Le rapport des
médianes d'une même paire varie jusqu'à un facteur 1,87 entre ordres.
Le solveur du quotient plus petit peut réduire le système dense de
taille $2^{s+r}$ à $2^s$, mais les produits de la série, la préparation
et les supports intermédiaires peuvent annuler cette économie.
Les entrées mesurées n'emploient qu'une ou deux directions non radicales
actives, même en grande dimension ambiante. Le smoke initial et les
écrans partagent encore la machine avec Étendue3D ; les diagnostics JIT
en processus partagé ne qualifient pas le temps froid. Aucun seuil
stable ni gain général ne découle de ces observations. Le manifeste
de 18\,000 identités, les quotients proches de la singularité et les
métriques générales restent à mesurer en conditions isolées.

\FloatBarrier
\subsection{Petit polynôme minimal}
\paragraph{Pourquoi ce choix ?}
Si une identité courte comme $A^2=\lambda 1$ est certifiée, une longue
puissance du \emph{même} multivecteur se réduit à deux coefficients
au lieu de reconstruire chaque puissance complète. La question
expérimentale est de savoir si le coût du certificat s'amortit lorsque
l'entrée et les requêtes se répètent.
Les inverses et fonctions de multivecteurs peuvent se déduire d'un
polynôme minimal, sans développer toutes les puissances dans la base
complète~\cite{prodanov}. Si $A=\sum_i P_i$, les $P_i$ anticommutent et
$P_i^2=s_i$ est scalaire, alors $A^2=\lambda=\sum_i s_i$ ; les puissances
et fonctions analytiques de $A$ restent dans $\operatorname{span}\{1,A\}$.
Cette identité promet un coût lié au degré minimal, une fois la structure
reconnue, pour des expressions répétées ou des inverses.

\Needspace{9\baselineskip}
\paragraph{Essai N04 : certifier puis réemployer.}
Le prototype N04 certifie rationnellement le carré complet
$A^2=\lambda 1$, lie le certificat aux coefficients et à la métrique,
et refuse une identité seulement approchée en flottant.
Le certificat est utilisable pour les puissances de la \emph{même} entrée :
un changement de coefficients ou de métrique l'invalide ; une entrée
refusée revient à une puissance complète exacte. Il ne s'agit pas de
chercher un polynôme minimal général. Sous Julia 1.13, les 1\,448/1\,448
assertions de correction et de contrat passent, ainsi que 70/70 tests du
contrôleur après correction d'un défaut de harnais.

Le premier smoke PerfChecker qualifie 20/20 cas, quatre fixtures et cinq
routes, avec 15 observations brutes par cas, soit 300 observations.
Chaque mesure est un épisode entier de longueur $H$ après préparation de
la fixture et de l'oracle, avec une sortie rationnelle complète et
possédée vérifiée à chaque étape. Les routes sont : puissance linéaire,
puissance binaire, DAG d'expression, certification à chaque demande,
et certification puis réemploi. Leurs coûts de calcul, de tentative de
certification, de repli et de sortie réellement exercés sont inclus.
Les quatre charges sont V2 : vecteur 2D stable, exposant 64, $H=1$ ;
M4 : entrée mixte admissible 4D stable, exposant 256, $H=32$ ;
N65 : métrique nulle 65D, entrée changeante, exposant 1024, $H=32$ ;
R4 : entrée 4D refusée par le certificat, exposant 16, $H=1$.

\begin{table}[ht]
\centering\small
\begin{tabular}{@{}llrrr@{}}
\toprule
Cas & Route & Médiane (\(\mu\)s) & Octets alloués & Allocations \\
\midrule
V2 & Linéaire & 80,939 & 254\,944 & 4\,790 \\
 & Binaire & 11,810 & 31\,504 & 573 \\
 & DAG & 29,643 & 51\,152 & 1\,099 \\
 & Certifier chaque fois & 8,455 & 12\,408 & 255 \\
 & Certifier puis réemployer & 7,857 & 12\,408 & 255 \\
\addlinespace
M4 & Linéaire & 56\,700,377 & 77\,051\,712 & 1\,707\,714 \\
 & Binaire & 509,630 & 1\,563\,680 & 30\,466 \\
 & DAG & 2\,997,477 & 4\,426\,800 & 102\,626 \\
 & Certifier chaque fois & 347,091 & 841\,008 & 18\,562 \\
 & Certifier puis réemployer & 169,484 & 375\,728 & 7\,433 \\
\addlinespace
N65 & Linéaire & 68\,078,844 & 12\,854\,272 & 147\,074 \\
 & Binaire & 3\,118,213 & 2\,879\,184 & 15\,234 \\
 & DAG & 12\,763,802 & 15\,234\,752 & 313\,570 \\
 & Certifier chaque fois & 12\,401,246 & 20\,226\,272 & 287\,746 \\
 & Certifier puis réemployer & 19\,261,904 & 29\,700\,984 & 419\,372 \\
\addlinespace
R4 & Linéaire & 67,183 & 201\,568 & 4\,184 \\
 & Binaire & 22,804 & 62\,608 & 1\,280 \\
 & DAG & 25,643 & 63\,248 & 1\,327 \\
 & Certifier chaque fois & 27,443 & 79\,152 & 1\,623 \\
 & Certifier puis réemployer & 29,486 & 79\,152 & 1\,623 \\
\bottomrule
\end{tabular}
\caption{Premier smoke N04 : médianes des 15 observations par cas,
en microsecondes pour l'épisode entier. Octets et nombres d'allocations
sont constants sur ces observations. Les quatre fixtures et les cinq
voies sont définies dans le texte.}
\label{tab:n04smoke}
\end{table}

\begin{figure}[ht]
\centering
\begin{tikzpicture}
\begin{semilogyaxis}[
 width=.91\linewidth,height=6.1cm,
 ybar,bar width=7pt,enlarge x limits=.17,
 ymin=5,ymax=100000,
 xtick={1,2,3,4},xticklabels={V2 stable,M4 stable,N65 changeant,R4 refus},
 ylabel={Médiane de l'épisode (\(\mu\)s)},
 grid=major,grid style={gray!20},
 legend style={at={(.5,1.03)},anchor=south,legend columns=3,
 font=\footnotesize,draw=none,fill=white},
]
\addplot+[fill=gray!55,draw=gray!75!black,mark=none] coordinates
 {(1,80.939) (2,56700.377) (3,68078.844) (4,67.183)};
\addplot+[fill=blue!55,draw=blue!80!black,mark=none] coordinates
 {(1,11.810) (2,509.630) (3,3118.213) (4,22.804)};
\addplot+[fill=orange!65,draw=orange!85!black,mark=none] coordinates
 {(1,29.643) (2,2997.477) (3,12763.802) (4,25.643)};
\addplot+[fill=green!55,draw=green!65!black,mark=none] coordinates
 {(1,8.455) (2,347.091) (3,12401.246) (4,27.443)};
\addplot+[fill=violet!55,draw=violet!70!black,mark=none] coordinates
 {(1,7.857) (2,169.484) (3,19261.904) (4,29.486)};
\legend{Linéaire,Binaire,DAG,Certifier chaque fois,Certifier puis réemployer}
\end{semilogyaxis}
\end{tikzpicture}
\caption{Les cinq routes N04 sur chaque fixture du
tableau~\ref{tab:n04smoke}. L'échelle logarithmique sert à voir ensemble
les épisodes courts et longs ; seules les routes d'un \emph{même} cas
sont comparables. Smoke à ordre fixe sous interférence Étendue3D.}
\label{fig:n04smoke}
\end{figure}

Sur V2, la certification réemployée est environ 1,5 fois plus rapide
que la puissance binaire. Sur M4, elle est environ 3 fois plus rapide
et alloue environ 4,2 fois moins d'octets. Sur N65, les coefficients
changent : la voie binaire est alors environ 6,2 fois plus rapide que
la certification réemployée et alloue environ 10,3 fois moins ; le
réemploi perd son intérêt dans cette fixture. Sur R4, la tentative
de certification refusée ajoute du travail avant le repli exact, et
la voie binaire a la médiane la plus basse. Ces contrastes décrivent
quatre fixtures, sans établir de seuil de routage.

La collecte s'est faite à un thread et dans un ordre fixe, sous charge
Étendue3D ; l'ordre inverse et un poste isolé restent à mesurer. La
préparation de l'algèbre et du plan varie entre les contrats des routes,
si bien que l'égalité des états préparés n'est pas encore garantie.
La grille prévue de 30\,720 identités sur douze dimensions n'a pas été
mesurée. Les cas $\lambda=0$, les changements d'entrée et les
perturbations figurent dans les contrôles, mais ne suffisent pas à
établir une politique générale.

\FloatBarrier

\subsection{Pfaffien pour une sortie scalaire}
\paragraph{Pourquoi ce choix ?}
Certaines applications demandent seulement le scalaire d'une longue
chaîne de vecteurs. Le calcul multivectoriel complet construit alors
des grades qui seront jetés. Le Pfaffien est testé pour éviter cette
expansion ; son coût de contractions et sa précision flottante doivent
être comparés au gain de travail.
L'expansion des produits de Clifford en contractions et Pfaffiens remonte
notamment à Wilmot~\cite{wilmotpf,wilmot2026}. Pour une forme bilinéaire
\emph{symétrique}, même singulière, sur un corps de caractéristique
différente de deux, la partie scalaire d'un produit \emph{ordonné} de
$2m$ vecteurs est une somme signée d'appariements. Avec
$K_{ij}=v_i\cdot v_j$ pour $i<j$ et $K_{ji}=-K_{ij}$,
\[
 \langle v_1\cdots v_{2m}\rangle_0=\operatorname{Pf}(K).
\]
La matrice $K$ est orientée par l'ordre des facteurs, et non obtenue par
antisymétrisation de la matrice de Gram symétrique. La relation de Clifford
donne la récurrence
\[
S(v_1,\ldots,v_{2m})=
\sum_{j=2}^{2m}(-1)^j(v_1\cdot v_j)
S(v_2,\ldots,\widehat v_j,\ldots,v_{2m}),
\]
avec $S(\varnothing)=1$ : c'est le développement du Pfaffien sur sa
première ligne, sans inversion de la métrique.
Ainsi, pour quatre vecteurs, la sortie vaut
$(v_1\!\cdot v_2)(v_3\!\cdot v_4)
-(v_1\!\cdot v_3)(v_2\!\cdot v_4)
+(v_1\!\cdot v_4)(v_2\!\cdot v_3)$.
Un algorithme dense de Pfaffien coûte cubiquement en $m$, après la
construction des contractions~\cite{wimmer}, alors qu'expanser le
multivecteur intermédiaire peut être exponentiel.
Pour $k=2m$ vecteurs de $n$ coordonnées, le prototype construit
densément les contractions en
$O(n^2k+nk^2)$ opérations avant son élimination $O(k^3)$ ;
ces coûts peuvent dominer quand seuls quelques axes sont actifs.

Cette voie ne calcule que le scalaire demandé, pas le produit complet.
Le prototype N05 vérifie les signes par un oracle rationnel indépendant
et refuse les sorties non scalaires. Son domaine,
les comparaisons chaudes et un contre-exemple de perte de précision
\code{Float64} suivent ce principe. Les
contractions et leur conditionnement font partie du coût et du contrat
numérique ; une identité rationnelle ne garantit pas une réponse
flottante exacte.
Pour plusieurs chaînes qui réutilisent un même pool de vecteurs, le
prolongement K3 partage les contractions et les buffers entre sorties
scalaires, avec vérification des changements de métrique et de
coordonnées. Il conserve la même élimination pfaffienne
et la même limite numérique. Le coût du plan, du workspace et des
contrôles décide si ce partage vaut la peine ; son essai suit celui de
N05.

\paragraph{Essai N05.}
\label{sec:n05perf}
Le prototype N05 calcule uniquement
$\langle v_1\cdots v_{2m}\rangle_0$, en construisant les contractions
puis en éliminant la matrice antisymétrique avec
pivotage. Cette élimination simple n'est pas l'implémentation
optimisée de Wimmer. Il passe 342/342 tests rationnels et de contrat sur
des métriques symétriques positives, signées, dégénérées, nulles et non
orthogonales, y compris en dimension 65. L'oracle indépendant construit
tous les grades par l'action extérieure de Chevalley ; il n'emploie ni
Pfaffien ni multiplication Garamon. Le smoke PerfChecker valide 15/15
routes, quatre fixtures et 225 observations à un thread, sous
interférence Étendue3D. Les quatre routes à contractions
\emph{préexistantes} appartiennent à une catégorie secondaire : leurs
temps ne sont pas comparés aux témoins qui reçoivent les vecteurs.

La figure~\ref{fig:n05temps} présente uniquement le contrat principal :
chaque épisode construit les contractions pour N05, ou les vecteurs et
le plan Garamon pour les témoins, puis conserve les scalaires de toutes
ses itérations. La voie récursive Garamon demande déjà la sortie
scalaire, mais n'admet ici que les métriques diagonales ; le produit
complet matérialise aussi les autres grades. Un témoin \emph{scalaire
ciblé} pour une métrique générale manque encore. Les durées couvrent
des longueurs et horizons distincts : la comparaison a un sens
\emph{à l'intérieur} de chaque fixture, pas entre les quatre groupes.

\begin{figure}[ht]
\centering
\begin{tikzpicture}
\begin{semilogyaxis}[
 width=.86\linewidth,height=5.8cm,
 ybar,bar width=7pt,enlarge x limits=.2,
 ymin=.2,ymax=3500,
 xtick={1,2,3,4},
 xticklabels={2D signée,65D signée,4D générale,8D dégénérée},
 x tick label style={rotate=15,anchor=east},
 ylabel={Médiane de l'épisode (\(\mu\)s)},
 grid=major,grid style={gray!20},
 legend style={at={(.5,1.03)},anchor=south,legend columns=3,column sep=10pt,draw=none,fill=white},
]
\addplot+[fill=gray!50,draw=black,mark=none] coordinates
 {(1,3.522) (2,1194.219) (3,40.469) (4,399.721)};
\addplot+[fill=orange!65,draw=orange!85!black,mark=none] coordinates
 {(1,11.565) (2,1729.413) (4,496.371)};
\addplot+[fill=blue!60,draw=blue!80!black,mark=none] coordinates
 {(1,.372) (2,342.305) (3,.595) (4,12.992)};
\legend{Complet,Récursif scalaire,Contractions + Pfaffien}
\end{semilogyaxis}
\end{tikzpicture}
\caption{Épisodes chauds du smoke N05 : longueurs de chaînes
$4,8,6,4$ et horizons $1,32,1,32$ dans l'ordre des groupes.
L'échelle est logarithmique. La barre récursive manque en 4D générale
par \emph{non-admission}, pas parce que son temps serait nul.
Construction des contractions incluse ; résultats exploratoires et
contrats de calcul différents pour le produit complet.}
\label{fig:n05temps}
\end{figure}

Le Pfaffien exploite ici la sortie restreinte et évite de produire tous
les grades, sans établir de seuil ni une supériorité générale. Son
identité algébrique exacte ne suffit pas à garantir la précision
flottante. Le contre-exemple exécuté
$G=\operatorname{diag}(1,1,-1)$,
$u=(1,2^{-30},1)$ donne exactement
$\langle u^4\rangle_0=2^{-120}$, mais la construction des contractions
en \code{Float64} suivie du Pfaffien retourne zéro : le petit terme
disparaît dans $1+2^{-60}-1$. Une tolérance absolue du smoke accepte ce
zéro et ne certifie donc pas les faibles scalaires. Une politique de
précision ou un repli exact, la construction spécialisée du Gram et une
reprise isolée précèdent toute admission au sélecteur. Le manifeste
étendu de 36\,480 identités, de 2 à 129 dimensions, est préparé
\emph{mais non exécuté}.

\FloatBarrier
\paragraph{Essai K3 : partager les contractions entre chaînes.}
\label{sec:k3perf}
Quand plusieurs chaînes réutilisent les mêmes vecteurs, K3 demande si
un plan et un workspace communs amortissent les contractions déjà
calculées, en tenant compte des contrôles et des changements d'entrées.
Le prolongement K3 de N05 compare cinq organisations du \emph{même}
calcul pfaffien scalaire : chaînes indépendantes, workspaces séparés,
partage avec allocations répétées, workspace commun et cache
vérifié. Chaque épisode reçoit une métrique, un
pool de vecteurs, plusieurs chaînes ordonnées et un horizon $H$,
puis rend une matrice \emph{possédée} de scalaires, une ligne par chaîne
et une colonne par itération. Le plan des paires, le workspace, la
construction des contractions, les contrôles de changement et les
copies de sorties sont payés dans l'épisode chaud ; aucun artefact
préparé n'est offert gratuitement.

Les 170 tests ciblés passent. Le smoke PerfChecker qualifie 20/20 routes,
300 observations, quatre fixtures et un thread sous interférence
Étendue3D. La figure~\ref{fig:k3partage} retient trois des cinq routes
pour montrer le contraste ; les deux autres, dont le partage avec
allocations répétées, appartiennent au même corpus. Les cinq voies
utilisent la même élimination et sont contrôlées contre l'oracle
rationnel, avec une tolérance \code{Float64} déclarée.

\begin{figure}[ht]
\centering
\begin{tikzpicture}
\begin{semilogyaxis}[
 width=.86\linewidth,height=5.8cm,
 ybar,bar width=7pt,enlarge x limits=.2,
 ymin=.1,ymax=3000,
 xtick={1,2,3,4},
 xticklabels={A : changeant,B : stable,C : une chaîne,D : disjoint},
 x tick label style={rotate=15,anchor=east},
 ylabel={Médiane de l'épisode (\(\mu\)s)},
 grid=major,grid style={gray!20},
 legend style={at={(.5,1.03)},anchor=south,legend columns=3,column sep=10pt,draw=none,fill=white},
]
\addplot+[fill=gray!50,draw=black,mark=none] coordinates
 {(1,29.708) (2,1425.417) (3,.229) (4,3.536)};
\addplot+[fill=blue!60,draw=blue!80!black,mark=none] coordinates
 {(1,61.315) (2,297.041) (3,29.993) (4,55.511)};
\addplot+[fill=orange!65,draw=orange!85!black,mark=none] coordinates
 {(1,72.010) (2,253.998) (3,30.725) (4,55.416)};
\legend{Indépendant,Workspace partagé,Cache vérifié}
\end{semilogyaxis}
\end{tikzpicture}
\caption{Coût total chaud de K3, préparation et sortie possédée
comprises, sur une échelle logarithmique. A : 4D, quatre chaînes,
$H=32$, coordonnées changeantes ; B : 65D, quatre chaînes, $H=32$,
coordonnées stables ; C : 4D, une chaîne, $H=1$ ; D : 4D générale,
quatre chaînes disjointes, $H=3$. Comparer seulement les routes
\emph{dans} une fixture. Ce smoke exploratoire ne fixe aucun seuil.}
\label{fig:k3partage}
\end{figure}

En B, les médianes sont 1\,425,417~\(\mu\)s pour les chaînes
indépendantes, 297,041 pour le workspace partagé et 253,998 pour le
cache vérifié. Un replay diagnostique séparé compte respectivement
128, 32 et 1 constructions de contractions ; ce compte n'est pas
ajouté au temps mesuré. En A, les changements imposent 32
reconstructions et le workspace partagé perd (61,315 contre
29,708~\(\mu\)s). En C, préparer le partage pour une seule courte
chaîne coûte près de 30~\(\mu\)s contre 0,229 en indépendant.
Moins de contractions ou d'allocations ne suffit donc pas à accélérer
l'épisode. Le partage ne corrige pas l'annulation \code{Float64} de
N05 : le contre-exemple $2^{-120}$ rendu nul passe aussi par ce
workspace.

\paragraph{Écran sélectif audité.}
Un essai distinct compare seulement les chaînes indépendantes au
workspace partagé, à sorties possédées, préparation et contrôles inclus.
Il qualifie 176/176 identités uniques sur onze dimensions, de 2 à 129,
avec horizons 1 et 32, entrées stables ou changeantes et deux ordres
de passage. L'audit vérifie les qualifications natives
et l'oracle rationnel sans manque ni problème de preuve. Cinq cas du
premier passage ont été remesurés lors de la reprise : les 2\,715
échantillons bruts restent conservés, mais ces doublons sont exclus
des 2\,640 échantillons et 176 cas retenus. Les 88 paires temporelles
comparent toujours deux routes \emph{dans le même passage}.

\begin{table}[ht]
\centering\small
\begin{tabular}{@{}lrrr@{}}
\toprule
Passage & Paires & Workspace plus rapide & Moins d'octets alloués \\
\midrule
Initial & 64 & 12 & 40 \\
Reprise & 24 & 20 & 24 \\
\bottomrule
\end{tabular}
\caption{K3, écran sélectif sous Julia 1.13. Les passages couvrent
des sous-ensembles de dimensions différents ; leurs comptes ne sont
pas fusionnés en une fréquence globale de gain. Un thread et
interférence Étendue3D déclarée.}
\label{tab:k3screen}
\end{table}

En 2D, $H=1$, entrée stable et ordre direct, l'épisode indépendant
prend 0,968~\(\mu\)s contre 52,146~\(\mu\)s pour le workspace.
En 129D, $H=32$, entrée stable et ordre direct, les médianes sont
4\,745,279 et 828,971~\(\mu\)s, avec respectivement
20\,339\,824 et 178\,592 octets alloués par épisode. Les chaînes
n'activent toutefois qu'au plus quatre directions et réutilisent
un pool limité : la dimension ambiante seule n'explique pas ces
contrastes. En 65D, $H=32$ et entrées changeantes, le rapport entre
routes varie avec l'ordre, surtout parce que le témoin indépendant
ralentit ; cette anomalie demande une reprise isolée, sans inférence
causale ni seuil.

Le premier passage a dépassé d'environ douze secondes sa fenêtre
nominale de terminaison ; la reprise propre s'est achevée sous son
plafond de 900 secondes. Sur les 440 identités de la grille sélective,
264 ne sont pas encore exécutées ; le manifeste distinct de
25\,920 identités de rejeu reste également non exécuté. Les cinq
routes, la précision flottante et les ordres devront être comparés
sur machine isolée avant tout choix automatique.

\FloatBarrier
\subsection{Rang de la matrice de Gram croisée}
\paragraph{Pourquoi ce choix ?}
Les produits de lames décomposables se décrivent par des contractions et
des angles entre sous-espaces~\cite{mandolesi}. Pour
$A=a_1\wedge\cdots\wedge a_p$ et
$B=b_1\wedge\cdots\wedge b_q$, formons
$C_{ij}=a_i\cdot b_j$. Les termes de contraction d'ordre $k$ font
intervenir des mineurs $k\times k$ de $C$ ; si son rang certifié est
$\rho$, ces mineurs sont nuls pour $k>\rho$. Une décomposition par
contractions pourrait donc omettre certains grades \emph{sans approximation}.

\paragraph{Expérience à conduire ; aucune mesure Garamon.}
Le filtre doit être prouvé pour la formule et la convention de produit
effectivement implémentées. Le test croisera rang faible et rang plein,
avec lames, métriques et sorties identiques ; il mesurera aussi le coût
de factorisation. Un seuil SVD flottant n'est qu'une décision approchée :
le chemin exact exige une certification de rang dans le corps choisi.

\subsection{États fermioniques gaussiens}
\paragraph{Pourquoi ce choix ?}
La représentation de fermions gaussiens conserve une description compacte
sous certaines transformations quadratiques~\cite{bravyigaussian}.
Elle suggère une représentation fermée pour des chaînes d'opérateurs de
Majorana ou de transformations apparentées, suivie d'une expansion seulement
à la sortie demandée. Les simulations non gaussiennes montrent qu'il faut
aussi suivre phases et écarts à la classe fermée~\cite{dias}.
En particulier, deux opérateurs ayant la même action par conjugaison sur
les vecteurs peuvent différer par un signe pertinent pour une amplitude.

\paragraph{Expérience à conduire ; aucune mesure Garamon.}
Le gain espéré concerne des chaînes longtemps fermées dans cette famille,
pas des multivecteurs arbitraires. L'essai vérifiera amplitudes et signes
avec un oracle complet sur de petits cas, puis introduira un nombre contrôlé
d'opérations non gaussiennes. Il enregistrera la croissance de la
représentation et le coût de la conversion finale.

\subsection{Changement de base factorisé}
\paragraph{Pourquoi ce choix ?}
Une matrice $Q$ sur les vecteurs induit $\Lambda^k Q$ sur le grade $k$,
matrice de taille $\binom nk\times\binom nk$. Des rotations orbitales
peuvent être appliquées comme une suite de rotations de Givens aux
coefficients concernés, sans matérialiser cette grande matrice~\cite{ffsim}.
Pour $Q$ orthogonale factorisée en $O(n^2)$ rotations, une application
directe coûte au plus de l'ordre de $n^2\binom nk$ opérations élémentaires
dans le cas dense ; la conversion et les accès mémoire doivent être comptés.

\paragraph{Expérience à conduire ; aucune mesure Garamon.}
La cible est le changement de base répété du générateur sur grades où la
matrice induite domine mémoire ou construction. Le benchmark comparera
application factorisée, matrice induite et calcul dans la base d'origine,
avec oracle des coefficients et mesure du pic mémoire. Le remplissage des
supports peut rendre la factorisation défavorable ; une transformation
non orthogonale demande une factorisation et un coût propres.

\subsection{Convolution disjointe signée}
\paragraph{Pourquoi ce choix ?}
Dans l'algèbre extérieure, le coefficient de masque $S$ du produit est
\[
 (a\wedge b)_S=
 \sum_{\substack{A\sqcup B=S}}(-1)^{\operatorname{inv}(A,B)}a_A b_B.
\]
Une sommation dense directe visite de l'ordre de $3^n$ couples disjoints.
Włodarczyk donne une convolution disjointe \emph{signée} en
$O^*(2^{\omega n/2})$ dans son modèle algébrique~\cite{wlodarczyk} ;
supprimer le signe ramènerait un autre produit. La borne asymptotique
n'indique pas les constantes, la mémoire ni le comportement sur supports
creux et coefficients flottants.

\paragraph{Expérience à conduire ; aucune mesure Garamon.}
Le test ne viserait que des produits extérieurs assez denses, à dimension
modérée, avec coefficients du domaine admis par l'algorithme. Il comparerait
temps et mémoire à la boucle creuse et au parcours complet, et certifierait
chaque coefficient. Les grandes tables auxiliaires et le coût des additions
peuvent annuler le gain asymptotique.

\subsection{Synthèse bilinéaire de petits noyaux}
\paragraph{Pourquoi ce choix ?}
Un produit fixé est un tenseur bilinéaire
$c_k=\sum_{ij}T_{kij}a_i b_j$. Une décomposition
$c=W((Ua)\odot(Vb))$ partage les multiplications entre sorties.
AlphaTensor a trouvé des décompositions de multiplication matricielle dans
un problème voisin~\cite{alphatensor}. Pour Garamon, l'intérêt serait une
synthèse \emph{hors ligne} de petits noyaux fixes, suivie d'une vérification
exacte du tenseur. Le nombre de multiplications seul ne prédit ni les
additions, ni la taille du code, ni la stabilité flottante.

\paragraph{Expérience à conduire ; aucune mesure Garamon.}
On comparera les noyaux synthétisés à la génération existante pour des
dimensions et grades bornés, en incluant génération, compilation et
exécution. Les identités devront être prouvées sur le corps de destination ;
une identité modulaire en caractéristique particulière ne s'étend pas
automatiquement à $\mathbb R$.

\subsection{Superoptimisation vérifiée}
\paragraph{Pourquoi ce choix ?}
Une fois le bon parcours choisi, un produit peut encore répéter des
millions de fois une courte séquence : calculer le signe d'une paire de
lames, leur XOR, leur intersection, puis l'adresse de sortie. C'est
\emph{ce coût local} qui motive la superoptimisation vérifiée. On cherche
une autre séquence d'instructions pour la même fonction, avec une preuve
d'équivalence pour les entiers et les cas limites admis. Elle ne change
ni les supports parcourus, ni les sorties demandées ; ainsi son intérêt
se mesurerait par appel du \emph{même} noyau, après avoir choisi sa
représentation et son algorithme.
Minotaur synthétise des séquences SIMD et vérifie leur équivalence
sémantique par Alive2~\cite{minotaur}; AlphaEvolve explore des programmes
par une autre méthode~\cite{alphaevolve}. Ces outils suggèrent de travailler
sur les petites boucles de signe, XOR, intersection et accumulation après
avoir fixé l'algorithme de haut niveau. Le gain recherché est local :
instructions, chargements et registres, non une réduction magique du
nombre de produits de lames.

\paragraph{Expérience à conduire ; aucun résultat de vitesse Garamon.}
Le témoin sera la boucle actuelle, face à une séquence candidate qui
reçoit exactement les mêmes masques, métriques et coefficients et renvoie
les mêmes sorties. On prouvera d'abord l'équivalence de la partie entière
pour les types, débordements et options de compilation utilisés. Pour les
accumulations flottantes, réassocier ou fusionner des opérations peut
changer les bits : ce contrat sera vérifié séparément. On mesurera
ensuite première compilation, temps chaud, allocations et taille du code
sur plusieurs architectures, avec la même stratégie de parcours des deux
côtés. Aucun candidat superoptimisé Garamon n'est encore validé ou
chronométré ; les figures de cet article ne donnent donc aucun gain pour
cette piste.

\subsection{Ordonnancement en fronts actifs}
\paragraph{Pourquoi ce choix ?}
Le \emph{wavefront path tracing} regroupe les chemins de calcul actifs
pour réduire la divergence sur GPU~\cite{lainewavefront}. Une récursion de
produit irrégulière pourrait pareillement produire des fronts de préfixes,
écarter les branches prouvées nulles, puis regrouper les survivantes par
grade ou destination. La réorganisation change l'ordonnancement et la
localité, pas la formule arithmétique.

\paragraph{Expérience à conduire ; aucune mesure Garamon.}
Le domaine plausible est un grand lot de produits partageant les mêmes
transitions. On comparera récursion directe et files de fronts sur CPU
puis, seulement si la granularité le justifie, sur GPU. La taille des files,
les transferts, la synchronisation et l'ordre des sommes flottantes sont
des coûts et des risques explicites ; une petite tâche creuse peut perdre.

\subsection{Index radix des supports matérialisés}
\paragraph{Pourquoi ce choix ?}
Les \emph{Adaptive Radix Trees} indexent des clés réellement stockées avec
une navigation compacte~\cite{art}. Ils pourraient accélérer les requêtes
sur un ensemble de masques \emph{matérialisé} et réutilisé. Cette idée est
distincte de l'arbre préfixe virtuel de la récursion C++ Garamon, où aucun
nœud d'index de support n'est nécessairement stocké.

\paragraph{Expérience à conduire ; aucune mesure Garamon.}
Le test comparera arbre radix, tableau trié, table de hachage et ZDD pour
recherches, intersections et itérations sur les mêmes supports ; il inclura
construction, mémoire retenue, localité et mises à jour. Un support aléatoire
peu partagé sert de cas négatif. L'index ne doit pas être introduit pour un
parcours unique qui ne l'amortit pas.

\Needspace{13\baselineskip}
\subsection{Arithmétique multimodulaire exacte}
\paragraph{Pourquoi ce choix ?}
Les tests d'exactitude eux-mêmes peuvent devenir coûteux : leur oracle
multiplie parfois de très grands entiers ou fractions pour contrôler tous
les coefficients. L'arithmétique multimodulaire a été retenue pour cette
\emph{arithmétique des coefficients}, distincte du rang binaire qui réduit
le nombre de \emph{masques} possibles et de la superoptimisation qui vise
les instructions d'une boucle. Elle calcule plusieurs petits résultats
modulaires, puis reconstruit le grand résultat exact. Elle pourrait
accélérer un oracle ou un produit entier répété si le coût des modules et
de la reconstruction reste inférieur au calcul direct en \code{BigInt} ;
elle ne vise pas les produits usuels \code{Float64}.
FLINT fournit réductions simultanées et reconstruction CRT avec
prétraitement réutilisable~\cite{flintcrt}. Pour des coefficients entiers ou
rationnels dont la taille rend l'arithmétique directe coûteuse, on pourrait
calculer modulo plusieurs premiers, puis reconstruire les coefficients.
Si le produit des modules dépasse deux fois une borne absolue certifiée,
la reconstruction signée donne l'entier exact. Un résidu non nul suffit à
prouver qu'un coefficient est non nul ; des résidus nuls sous quelques
premiers ne prouvent pas qu'il est zéro.

\paragraph{Expérience à conduire ; aucun résultat de vitesse Garamon.}
Le témoin calculera directement les mêmes oracles symboliques, inverses
et produits à grands coefficients en \code{BigInt} ou en rationnels.
La voie candidate paiera chaque réduction modulo un premier, chaque
produit modulaire et la reconstruction CRT. Une borne certifiée sur les
coefficients, les dénominateurs et les mauvais premiers conditionnera
l'égalité exacte de chaque résultat reconstruit. On comparera temps
\emph{total} et mémoire à entrées identiques, sur plusieurs tailles de
coefficients pour voir à partir de quand la préparation s'amortit.
Aucun prototype multimodulaire de Garamon n'a encore été mesuré ; les
temps des autres figures ne démontrent donc rien pour cette voie.

\subsection{Prédicats géométriques filtrés}
\paragraph{Pourquoi ce choix ?}
CGAL associe approximation, borne d'erreur et repli exact quand le signe
reste incertain~\cite{cgalfiltered}. Une API Garamon ne demandant
que le signe d'une orientation, d'une incidence ou d'une comparaison
pourrait suivre ce schéma : retourner vite si l'intervalle
exclut zéro, sinon calculer la décision exacte. La correction porte sur
le prédicat relativement à l'encodage exact de ses entrées.

\paragraph{Expérience à conduire ; aucune mesure Garamon.}
Le test séparera décisions loin de zéro et cas presque dégénérés, en
enregistrant le taux de repli et son coût. Les constructions intermédiaires
inexactes et l'origine des coordonnées doivent être spécifiées. Un
prédicat exact ne remplace ni une sortie multivectorielle complète ni
un calcul approché assorti d'une simple tolérance.

\paragraph{Admission structurelle.}
Ces voies supposent des certificats attachés aux objets ou aux plans :
rang binaire, secteur, radical, polynôme minimal, décomposabilité, rang
de Gram ou borne de coefficients. Ils sont établis une fois, propagés
quand la règle est prouvée et invalidés lors d'une opération non couverte.
La borne $2^d$ illustre pourquoi la réduction peut être majeure, mais
chaque coût de certification et chaque conversion doivent
figurer dans les mesures. Aucun verdict de compatibilité avec le contrat
public de Garamon n'est arrêté ici. Une matrice séparée inventorie
46 familles et 1\,035 paires avec leurs conditions et tests
proposés ; elle ne contient aucun temps nouveau et
ne démontre aucun gain composé.
\figureattente{Voies structurelles et compositions exactes}
{Abscisse : famille de structure et taille active ; ordonnées : coût de
certification, temps total, mémoire maximale et taux de refus.}
{Synthèse prévue des pistes sans résultat de vitesse qualifié ; chaque
comparaison devra partager la même sortie et son oracle.}
{fig:voies-structurelles}

\section{Discussion et pistes expérimentales}
Les méthodes ne forment pas une hiérarchie unique. Les travaux sur la
planification de programmes creux~\cite{galley}, les workspaces
\cite{sparseworkspace}, les jointures~\cite{leapfrog} et l'égalité saturée
\cite{egg} fournissent des idées transverses ; les intégrer à Garamon exige
des preuves de règles et des mesures sur ses propres charges. Les hypothèses
suivantes restent ouvertes :
\begin{enumerate}
\item Comparer accumulateur haché, tampon dense à positions touchées, fusion
triée, blocs de grades et composition de quatre facteurs ou plus. Mesurer
préparation, conversion, JIT, exécution et mémoire maximale, y compris en
chaînes longues et avec supports changeants.
\item Séparer prétraitement, première compilation, cache natif entre processus,
code conservé, plans et workspaces. Précompiler éventuellement un catalogue
\emph{fini} d'algèbres courantes, puis tester aussi la découverte imprévisible
d'une métrique ou d'un support. Le coût de construction du cache ou d'une image
applicative doit être amorti sur les sessions réellement visées.
\item Comparer un thread, plusieurs threads Julia et plusieurs processus sur
produit isolé et lots indépendants. Partager les plans immuables, allouer un
workspace par tâche, compter démarrage, sérialisation, duplication de mémoire
et éventuelle contention. Sur GPU, mesurer également transfert, lancement,
synchronisation et VRAM. Aucun gain GPU n'est encore établi.
\item Étendre et qualifier le premier portefeuille contextuel entre calcul
complet, récursion,
grades, jointure, plan, workspace, génération et représentation structurée.
Comparer règle analytique, table de coûts, puis modèle appris avec optimisation
des hyperparamètres. Séparer familles d'algèbres et patrons de supports entre
apprentissage et test. Les valeurs de Shapley/SHAP~\cite{shap} pourraient
\emph{expliquer} un modèle retenu, sans prouver sa causalité ni sa performance.
Mesurer regret devant le meilleur choix rétrospectif, surcoût du sélecteur et
qualité du premier appel.
\end{enumerate}

Une expérience d'élagage \emph{approché} peut omettre une contribution, la
journaliser et la corriger plus tard. Pour une récurrence affine, si
$x_{t+1}=L_tx_t+b_t$ et
$\widehat{x}_{t+1}=L_t\widehat{x}_t+b_t-d_t$, l'erreur
$e_t=x_t-\widehat{x}_t$ suit
\begin{equation}
e_{t+1}=L_t e_t+d_t.
\label{eq:recurrence}
\end{equation}
Le seul fait de posséder une récurrence ne reconstruit donc pas les termes
oubliés : il faut transporter leur résidu, borner son influence, ou rejouer un
segment depuis un état exact. Pour une récurrence non linéaire, les termes
croisés compliquent encore la correction. Une roulette sur une série de
corrections télescopiques peut être non biaisée sous conditions~\cite{rhee}, sans garantir une
trajectoire exacte à chaque réalisation. Les futures comparaisons distingueront
égalité rationnelle, résultat flottant bit à bit, erreur certifiée et simple
erreur observée. Le coût du certificat, du journal et du replay appartient au
temps total. Une variante approximative ne remplace jamais le contrat exact
sans demande explicite de l'utilisateur.

\section{Limites de validité et reproductibilité}
Les campagnes couvrent des domaines et budgets distincts. Un support fixe
peut rester calculable quand $n$ croît ; aucun « dernier $n$ » universel
n'existe. Le balayage entier jusqu'à 256 dimensions ne couvre pas toutes les
métriques ni les expressions longues. Les mesures de workspace portent
principalement sur une métrique euclidienne et sur des résultats consommés
immédiatement. Les temps de 10\,000 appels n'ont qu'un échantillon par cas.
La comparaison \code{UInt128}/\code{BigInt} isole un noyau simplifié.
L'expérience de pruning conserve au plus quatre directions actives et ne
représente pas le noyau public. L'étude du catalogue natif mesure la réutilisation du
cache natif après redémarrage pour un seul catalogue EGA3 et douze contextes
4D ; elle ne qualifie pas un catalogue de production ni une autre machine.
La première campagne TripleJoin n'utilise que deux à quatre directions actives,
et les longues traces n'ont parfois qu'un échantillon. Des processus
Étendue3D concurrents ont pu modifier la contention et la fréquence du CPU
durant cette phase expérimentale. Tous les classements de vitesse de cette
version doivent donc être relus comme un criblage ; une répétition
\emph{intégrale et isolée} est prévue avant d'arrêter des seuils.

N01 a été évalué à un thread et en deux ordres sous cette même
interférence ; les 648 cas sont des épisodes directement chronométrés,
mais les 10\,287 cas de son manifeste étendu ne sont pas encore exécutés.
Le rejeu distinct K1v2 n'a qualifié que son premier lot de dimension 2,
soit 324/41\,148 cas ; les autres lots demeurent non exécutés.

N03 dispose désormais de deux écrans corrigés de 756 cas chacun,
sur douze dimensions et en ordre inversé. Leur interférence machine,
leur faible support actif et les choix sensibles à l'ordre empêchent
encore tout seuil de vitesse ; les diagnostics JIT en processus partagé
ne qualifient pas le froid. Le premier essai rejeté pour erreur du
générateur reste archivé à part, hors des 1\,512 cas retenus.

N05 ne couvre que quatre fixtures de chaînes de vecteurs et une sortie
scalaire ; la construction des contractions est chronométrée, mais le
témoin scalaire ciblé en métrique générale et la politique de précision
flottante manquent. Ses 36\,480 identités de rejeu ne sont pas exécutées.

K3 possède quatre fixtures de smoke à cinq routes, puis un écran
sélectif audité de 176 identités uniques sur onze dimensions et deux
routes. Le premier passage a dépassé sa fenêtre nominale d'environ
douze secondes ; les cinq cas remesurés sont exclus des comparaisons.
Les 264 identités sélectives restantes, les 25\,920 identités du
rejeu complet, les trois autres routes et la validation numérique sur
des entrées mal conditionnées restent ouvertes.
N04 passe ses assertions rationnelles sous Julia 1.13 et son premier
smoke qualifie 20/20 cas avec 300 observations, sous interférence
Étendue3D et dans un seul ordre. Les préparations diffèrent entre
routes ; sa grille de 30\,720 identités n'est pas encore mesurée.
Le troisième lot B1 qualifie 96 épisodes paramétrés, de 1 à 1\,024
produits, avec 2\,016 observations sous la même interférence. Ses
supports ont seulement quatre masques et ses résultats ne qualifient
ni la piste B1W ni un choix général de contrôles de bornes.

Tant que des traitements Étendue3D partagent la machine, les prochains
essais de criblage seront limités à quatre threads ou workers au plus,
le plus souvent un. La grande matrice de cas sera rejouée ensemble
dans un environnement stable, sans autres jobs concurrents, avec les
mêmes contrats et des observations conservées.

Le banc C++/Julia apparié ne couvre que le produit packed dans deux algèbres,
deux supports et quatre horizons. Ses 31 échantillons individuels n'ont pas
été conservés dans la première passe ; seules leurs statistiques et les
résultats d'oracle restent disponibles. Il ne faut pas interpréter
\code{cpp\_upstream} comme une quatrième modalité appariée.
\section{Contrôles de bornes et coût de la validation}
\paragraph{Options de compilation et contrôles de bornes.}
Le noyau C++ packed apparié est compilé avec Clang, \code{-O2},
\code{-march=native} et \code{-ffp-contract=off} ; GCC n'est pas une
modalité de cet essai. La boucle Julia 1.13 \code{run\_packed\_batch} porte déjà
\code{@inbounds}. L'écart observé entre ces deux noyaux ne peut donc
pas être attribué, sur cette seule base, à la suppression des contrôles
de bornes ou à un hypothétique essai GCC.

Le pilote noté B1, séparé du noyau public,
compare une accumulation bufferisée vérifiée et une variante annotée.
À chaque appel, il copie les masques et chemins dans des vecteurs privés,
valide algèbre, supports, tailles et bornes de tous les indices ; seule
l'accumulation sur les buffers privés utilise \code{@inbounds}. Les deux
variantes paient les mêmes copies et validations. Sous Julia 1.13, les
tests réussissent 550/550 avec \code{--check-bounds=yes}, puis 550/550
avec \code{--check-bounds=auto}.

\paragraph{Premier smoke B1 : code spécialisé et temps complet.}
Sous Julia 1.13, un premier essai PerfChecker qualifie 8/8 cas en 8D
euclidienne à petit support et $H=1$. Deux processus distincts emploient
\code{--check-bounds=yes} ou \code{auto}, avec \code{-O2},
\code{--cpu-target=native} et un thread. Chaque processus passe les voies
vérifiée et annotée dans les deux ordres. Les trois phases --- construction
du plan, calcul chaud à plan réutilisé, et épisode complet avec plan
reconstruit et sortie possédée --- ont chacune sept observations par cas :
168 observations au total. L'épisode est chronométré directement et
l'oracle entier indépendant vérifie les sorties.

\begin{table}[ht]
\centering\small
\begin{tabular}{@{}llrr@{}}
\toprule
Bornes & Voie & Ordre 1 : chaud / épisode & Ordre 2 : chaud / épisode \\
\midrule
\code{yes} & Vérifiée & 1,122 / 6,361 & 0,867 / 6,256 \\
\code{yes} & Annotée & 0,977 / 6,442 & 0,901 / 6,320 \\
\code{auto} & Vérifiée & 0,887 / 7,331 & 0,941 / 6,120 \\
\code{auto} & Annotée & 0,840 / 6,915 & 0,857 / 6,263 \\
\bottomrule
\end{tabular}
\caption{Premier smoke B1 : médianes en microsecondes de sept
observations par phase et par cas. Le calcul chaud réemploie le plan ;
l'épisode inclut sa reconstruction. Les deux ordres sont des passages
distincts sous interférence Étendue3D.}
\label{tab:b1smoke}
\end{table}

Toutes les voies allouent les mêmes 3\,536 octets en 42 allocations dans la
phase chaude, puis 12\,640 octets en 221 allocations dans l'épisode ; la
construction seule alloue 9\,104 octets en 179 allocations. Hors mesure,
l'IR LLVM spécialisé contient un chemin \code{throw\_boundserror} dans
l'accumulation vérifiée sous \code{auto}, absent de l'accumulation
annotée ; sous \code{yes}, ce chemin apparaît dans les deux. La
différence de code recherchée existe donc pour cette boucle, sans prouver
que tous les contrôles du produit complet ont disparu. Les écarts chauds
ne donnent pas de gain stable sur l'épisode avec préparation ; l'ordre,
les sept observations par phase et la charge Étendue3D empêchent de
conclure à un réglage de production ou à un gain général.

\paragraph{Prévol B1 : douze dimensions ambiantes.}
Le premier smoke ne couvre que 8D. Un balayage suivant demande si son
contraste chaud résiste à une variation de la dimension ambiante et au
changement de représentation des masques après 64D. Il qualifie 48/48 cas
sur les douze dimensions du tableau~\ref{tab:b1dimensions}, sous Julia 1.13 :
deux voies et deux ordres par dimension. Chaque dimension a son processus
neuf, avec \code{--check-bounds=auto}, \code{-O2}, cible \code{native},
un thread, métrique positive, support \code{low\_grade} et $H=1$.
Construction, calcul chaud et épisode complet fournissent sept
observations par cas et par phase, soit 1\,008 observations brutes ;
l'oracle exact et la propriété des sorties sont vérifiés.

\begin{table}[ht]
\centering\footnotesize
\begin{tabular}{@{}rrrrrrr@{}}
\toprule
$n$ & Chaud V & Chaud A & Épisode V & Épisode A & Octets chaud & Octets épisode \\
\midrule
2 & 0,810 & 0,828 & 9,470 & 6,285 & 3\,536 & 12\,496 \\
3 & 0,831 & 0,796 & 6,151 & 6,162 & 3\,536 & 12\,496 \\
4 & 1,008 & 0,821 & 6,094 & 5,989 & 3\,536 & 12\,544 \\
6 & 0,860 & 0,802 & 6,627 & 6,359 & 3\,536 & 12\,592 \\
8 & 0,864 & 0,824 & 6,520 & 6,533 & 3\,536 & 12\,640 \\
12 & 0,869 & 0,825 & 6,347 & 6,483 & 3\,536 & 12\,736 \\
16 & 0,879 & 0,819 & 6,378 & 6,341 & 3\,536 & 12\,832 \\
32 & 0,888 & 0,822 & 6,629 & 6,311 & 3\,536 & 13\,216 \\
64 & 0,972 & 0,921 & 6,569 & 6,393 & 3\,536 & 13\,984 \\
65 & 1,125 & 1,175 & 7,316 & 7,172 & 4\,384 & 17\,312 \\
96 & 1,124 & 1,150 & 7,550 & 7,185 & 4\,384 & 18\,128 \\
128 & 1,129 & 1,051 & 7,280 & 7,617 & 4\,384 & 18\,944 \\
\bottomrule
\end{tabular}
\caption{Prévol B1 : médianes en microsecondes sur les deux passages,
soit 14 observations par voie et phase. V est la voie vérifiée, A la
voie annotée ; l'épisode reconstruit le plan et conserve la sortie.
Les octets alloués sont identiques entre V et A.}
\label{tab:b1dimensions}
\end{table}

\begin{figure}[ht]
\centering
\begin{tikzpicture}
\begin{axis}[
 width=.86\linewidth,height=5.3cm,
 xmin=.7,xmax=12.3,ymin=0,ymax=20500,
 xtick={1,2,3,4,5,6,7,8,9,10,11,12},
 xticklabels={2,3,4,6,8,12,16,32,64,65,96,128},
 xlabel={Dimension ambiante $n$},ylabel={Octets alloués},
 grid=major,grid style={gray!20},
 legend style={at={(.03,.97)},anchor=north west,draw=none,fill=white},
]
\addplot+[blue!75!black,thick,mark=*] coordinates
 {(1,3536) (2,3536) (3,3536) (4,3536) (5,3536) (6,3536)
  (7,3536) (8,3536) (9,3536) (10,4384) (11,4384) (12,4384)};
\addplot+[orange!85!black,thick,mark=square*] coordinates
 {(1,12496) (2,12496) (3,12544) (4,12592) (5,12640) (6,12736)
  (7,12832) (8,13216) (9,13984) (10,17312) (11,18128) (12,18944)};
\draw[dashed,gray!70!black] (axis cs:9.5,0) -- (axis cs:9.5,20500);
\legend{Calcul chaud,Épisode complet}
\end{axis}
\end{tikzpicture}
\caption{Allocations B1 des deux voies, identiques sur chaque cas.
La ligne pointillée sépare 64D et 65D, où la fixture passe de
\code{UInt64} à \code{UInt128}. Les masques d'entrée restent
$0,1,2,3$ : ce graphique ne mesure pas la croissance d'un support
algébrique plus dense.}
\label{fig:b1dimensions}
\end{figure}
\FloatBarrier

Le calcul chaud effectue 42 allocations à chaque dimension et par voie ;
l'épisode en effectue 221 de 2 à 64D et 273 dès 65D. La rupture des
octets et des comptes à 65D coïncide avec le passage des masques de
\code{UInt64} à \code{UInt128}. Les quatre masques d'entrée sont toujours
$0,1,2,3$ : seules deux directions de base sont actives, quelle que soit
la dimension ambiante. Celle-ci change surtout les métadonnées de
l'algèbre et le type de masque ; les grades élevés et supports denses
restent hors de ce prévol.

Les temps chauds favorisent souvent légèrement l'annotation de 3 à 64D,
mais l'épisode complet ne montre aucun gain stable. En 2D, l'épisode
vérifié du second passage a une médiane de 10,874~\(\mu\)s contre
6,149~\(\mu\)s au premier, alors que ses temps chauds sont
0,811 et 0,809~\(\mu\)s : la médiane regroupée de 9,470~\(\mu\)s
n'isole donc pas un effet causal des bornes. En 8D les épisodes sont
presque égaux et en 128D la voie vérifiée a la médiane la plus basse.
Sous interférence Étendue3D, ces 48 cas ne fixent ni seuil, ni option
par défaut, ni comparaison C++.

\paragraph{Prévol B1 : lames de grade élevé.}
Un second lot éprouve la limite du support précédent : les quatre
masques d'entrée deviennent $0,1,2^n-1,2^n-2$. Les deux derniers ont
les grades $n$ et $n-1$, tandis que le nombre de termes par opérande et
de chemins reste borné. Les douze dimensions et les options sont les
mêmes que pour le bas grade : \code{auto}, \code{-O2}, cible
\code{native}, métrique positive, $H=1$, un thread et deux ordres.
Ce lot qualifie aussi 48/48 cas, avec sept observations pour chacune
des trois phases par cas, soit 1\,008 observations brutes ; l'oracle
exact et les sorties possédées sont vérifiés. Les deux prévols dimensionnels
B1 totalisent ainsi 96 cas et 2\,016 observations.

\begin{table}[ht]
\centering\footnotesize
\begin{tabular}{@{}rrrrrr@{}}
\toprule
$n$ & Chaud V & Chaud A & Épisode V & Épisode A & Octets épisode \\
\midrule
2 & 0,825 & 0,760 & 6,173 & 6,463 & 12\,496 \\
3 & 1,123 & 0,784 & 6,883 & 6,602 & 12\,496 \\
4 & 0,792 & 0,840 & 6,265 & 6,096 & 13\,024 \\
6 & 0,770 & 0,778 & 6,796 & 6,389 & 13\,744 \\
8 & 0,762 & 0,803 & 6,306 & 6,611 & 13\,792 \\
12 & 0,786 & 0,762 & 8,101 & 7,003 & 17\,888 \\
16 & 0,810 & 0,790 & 6,893 & 7,540 & 17\,984 \\
32 & 0,849 & 0,812 & 7,133 & 6,907 & 18\,368 \\
64 & 0,890 & 0,885 & 8,359 & 8,220 & 33\,152 \\
65 & 1,153 & 1,240 & 10,611 & 10,300 & 36\,128 \\
96 & 1,191 & 1,160 & 11,265 & 11,591 & 36\,944 \\
128 & 1,260 & 1,196 & 12,735 & 12,160 & 37\,760 \\
\bottomrule
\end{tabular}
\caption{Prévol B1 sur lames de hauts grades : médianes en
microsecondes sur 14 observations par voie et phase, soit deux passages
de sept. V est la voie vérifiée et A la voie annotée. L'épisode inclut
la construction du plan et une sortie possédée ; ses octets alloués sont
identiques entre V et A.}
\label{tab:b1highgrade}
\end{table}

\begin{figure}[ht]
\centering
\begin{tikzpicture}
\begin{axis}[
 width=.86\linewidth,height=5.3cm,
 xmin=.7,xmax=12.3,ymin=0,ymax=40500,
 xtick={1,2,3,4,5,6,7,8,9,10,11,12},
 xticklabels={2,3,4,6,8,12,16,32,64,65,96,128},
 xlabel={Dimension ambiante $n$},ylabel={Octets alloués par épisode},
 grid=major,grid style={gray!20},
 legend style={at={(.03,.97)},anchor=north west,draw=none,fill=white},
]
\addplot+[blue!75!black,thick,mark=*] coordinates
 {(1,12496) (2,12496) (3,12544) (4,12592) (5,12640) (6,12736)
  (7,12832) (8,13216) (9,13984) (10,17312) (11,18128) (12,18944)};
\addplot+[orange!85!black,thick,mark=square*] coordinates
 {(1,12496) (2,12496) (3,13024) (4,13744) (5,13792) (6,17888)
  (7,17984) (8,18368) (9,33152) (10,36128) (11,36944) (12,37760)};
\legend{Bas grades,Grades $n$ et $n-1$}
\end{axis}
\end{tikzpicture}
\caption{Allocations des épisodes complets B1 dans les deux lots.
Les voies vérifiée et annotée ont les mêmes allocations à dimension
et support fixés. Les lots ont été exécutés successivement sous
Étendue3D : cette différence descriptive ne suffit pas à attribuer
causalement tout le surcoût aux grades.}
\label{fig:b1gradealloc}
\end{figure}
\FloatBarrier

Le calcul chaud alloue 3\,536 octets en 42 allocations jusqu'à 64D,
puis 4\,384 octets en 42 allocations dès 65D, comme dans le lot de
bas grade. L'épisode de hauts grades passe de 221 allocations en 2--3D
à 279 en 64D et 309 dès 65D ; ses octets augmentent surtout avec
la construction du plan, tandis que l'accumulation à quatre termes
reste courte. À 8D, l'épisode vérifié est plus court dans cette archive ;
à 12D, l'annoté l'est, puis le signe change encore à 16D et à 96D.
Aucun avantage d'épisode uniforme ne se dégage. Les deux lots ont été
exécutés successivement sous interférence Étendue3D : leurs options
sont identiques, mais leur écart de temps ne mesure pas strictement
le seul effet du grade. Le profilage de la préparation et un rejeu
isolé précèdent toute règle de sélection.
\FloatBarrier

\paragraph{Prévol B1 : épisodes de 1 à 1\,024 produits.}
\emph{Pourquoi ce troisième essai ?} Les deux lots précédents mesurent
$H=1$ : ils ne disent pas si la petite différence du calcul chaud
s'accumule lorsque les mêmes supports servent longtemps. Ce lot compare
donc les deux voies sur des épisodes complets de $H=1,32,1\,024$, en
comptant la construction du plan une fois, les validations et copies
privées à \emph{chaque} produit, et la conservation de toutes les sorties.
Il croise 2, 8, 65 et 128D avec les supports de bas et de hauts grades,
deux voies et deux ordres, sous Julia 1.13, \code{--check-bounds=auto},
\code{-O2}, cible \code{native}, métrique positive et un thread. Les
96/96 cas sont qualifiés ; sept observations par phase et par cas donnent
2\,016 observations brutes. L'oracle exact vérifie les $H$ sorties
possédées et indépendantes avant et après la mesure. La charge Étendue3D
rend ces temps exploratoires.

\begin{table}[ht]
\centering\small
\begin{tabular}{@{}rlrrrrr@{}}
\toprule
$n$ & Support & $H$ & V (\(\mu\)s) & A (\(\mu\)s) & Octets & Alloc. \\
\midrule
2 & bas & 1 & 6,395 & 6,242 & 12\,496 & 221 \\
2 & bas & 32 & 24,383 & 23,586 & 114\,736 & 1\,399 \\
2 & bas & 1\,024 & 602,234 & 580,655 & 3\,384\,312 & 39\,096 \\
2 & élevé & 1 & 6,862 & 6,796 & 12\,496 & 221 \\
2 & élevé & 32 & 24,969 & 23,578 & 114\,736 & 1\,399 \\
2 & élevé & 1\,024 & 585,426 & 581,638 & 3\,384\,312 & 39\,096 \\
8 & bas & 1 & 6,797 & 6,444 & 12\,640 & 221 \\
8 & bas & 32 & 24,183 & 25,011 & 114\,880 & 1\,399 \\
8 & bas & 1\,024 & 628,263 & 613,279 & 3\,384\,456 & 39\,096 \\
8 & élevé & 1 & 7,158 & 6,698 & 13\,792 & 233 \\
8 & élevé & 32 & 25,785 & 25,227 & 116\,032 & 1\,411 \\
8 & élevé & 1\,024 & 624,698 & 631,071 & 3\,385\,608 & 39\,108 \\
65 & bas & 1 & 7,223 & 7,037 & 17\,312 & 273 \\
65 & bas & 32 & 32,215 & 33,181 & 145\,840 & 1\,451 \\
65 & bas & 1\,024 & 895,384 & 886,142 & 4\,256\,632 & 39\,148 \\
65 & élevé & 1 & 10,235 & 10,722 & 36\,128 & 309 \\
65 & élevé & 32 & 40,293 & 40,668 & 164\,656 & 1\,487 \\
65 & élevé & 1\,024 & 1\,028,550 & 1\,029,182 & 4\,275\,448 & 39\,184 \\
128 & bas & 1 & 7,283 & 7,107 & 18\,944 & 273 \\
128 & bas & 32 & 33,577 & 35,022 & 147\,472 & 1\,451 \\
128 & bas & 1\,024 & 890,260 & 912,010 & 4\,258\,264 & 39\,148 \\
128 & élevé & 1 & 12,680 & 12,262 & 37\,760 & 309 \\
128 & élevé & 32 & 43,479 & 42,829 & 166\,288 & 1\,487 \\
128 & élevé & 1\,024 & 1\,074,436 & 1\,049,505 & 4\,277\,080 & 39\,184 \\
\bottomrule
\end{tabular}
\caption{Troisième prévol B1 : médianes d'épisode complet sur deux
ordres, 14 observations par voie et cellule. V est la voie vérifiée,
A la voie annotée. Octets et allocations sont identiques entre voies
à cellule égale ; toutes les sorties sont possédées. Les temps des
lots successifs ne constituent pas une comparaison appariée entre
grades.}
\label{tab:b1horizons}
\end{table}

\begin{figure}[!b]
\centering
\begin{tikzpicture}
\begin{axis}[
 width=.86\linewidth,height=5.1cm,
 ybar,bar width=13pt,enlarge x limits=.08,
 ymin=-3,ymax=5,ytick={-3,-2,-1,0,1,2,3,4,5},
 xtick={1,2,3,4,5,6,7,8},
 xticklabels={2 bas,2 haut,8 bas,8 haut,65 bas,65 haut,128 bas,128 haut},
 x tick label style={rotate=35,anchor=east},
 ylabel={Avantage de A sur V (\%)},
 grid=major,grid style={gray!20},
]
\addplot+[fill=blue!60,draw=blue!80!black,mark=none] coordinates
 {(1,3.583) (2,.647) (3,2.385) (4,-1.020)
  (5,1.032) (6,-.061) (7,-2.443) (8,2.320)};
\end{axis}
\end{tikzpicture}
\caption{Écart relatif des médianes d'épisode à $H=1\,024$,
$100(\mathrm{V}-\mathrm{A})/\mathrm{V}$ : une barre positive favorise
l'annotation. Les écarts restent petits dans ce prévol sous Étendue3D ;
la barre 65D haut grade est pratiquement nulle. Ce graphique ne mesure
pas le gain potentiel de B1W.}
\label{fig:b1horizons}
\end{figure}
\FloatBarrier

À $H=1\,024$, le plus grand gain descriptif de l'annotation est de
3,6\% en 2D bas grade. Deux reculs dépassent 1\% et un troisième
écart défavorable de 0,06\% est une quasi-égalité : la médiane annotée
est supérieure dans trois des huit cellules. Les voies ont exactement
les mêmes allocations dans chaque cellule. La différence de code LLVM
dans la boucle ne devient donc pas un gain d'épisode uniforme lorsque
copies, validations, lectures et sorties possédées sont répétées.
L'interférence et les écarts proches interdisent un choix de production
fondé sur ces seules médianes.

\paragraph{Hypothèse B1W : réemployer un instantané validé.}
Les 39\,096 à 39\,184 allocations et 3,38 à 4,28 millions d'octets
alloués à $H=1\,024$ orientent une autre expérience. B1W conserverait
un instantané validé du plan ainsi que ses chemins et buffers privés
pendant l'épisode. Les coefficients changeants seraient relus à chaque
produit ; les $H$ sorties resteraient exactes, possédées et
indépendantes. Le prototype devra définir et tester l'invalidation lors
d'une mutation du plan, de l'algèbre ou des supports, puis mesurer la
validation initiale, la mémoire retenue, les invalidations, le temps et
les allocations d'épisodes complets face aux voies B1 actuelles. Les
copies nécessaires aux sorties possédées ne pourront pas disparaître.
B1W n'est pas encore implémenté ni chronométré : aucun gain ne lui est
attribué. La reprise isolée devra également croiser métriques mixtes ou
dégénérées, supports centraux et plus denses, \code{-O3} et cible
\code{generic} avant tout transfert vers le noyau public.
\FloatBarrier

Le patch C++ de déplacement a son \emph{propre} témoin amont, avec mêmes
entrées et même représentation \code{Mvec}. Ses compteurs d'allocations
sont déterministes sur EGA3/EGA7 euclidiennes en \code{Float64} ; ses temps
restent exploratoires sous la même activité concurrente. Les résultats
d'allocations ne démontrent pas un gain sur d'autres algèbres, scalaires
ou compilateurs.

\Needspace{10\baselineskip}
\section{Transfert et programme de validation}
\subsection{Ordre des prochaines validations}
Le programme expérimental suit quatre étapes, dans cet ordre :
\begin{enumerate}
\item \textbf{Prévol et validation des techniques.} Définir pour chaque
technique son contrat de sortie, ses domaines admissibles, ses budgets et
ses témoins ; vérifier la correction avec un oracle indépendant, les
refus attendus et la reproductibilité des identités sous Julia 1.13.
C'est l'étape active pour les prototypes encore partiellement couverts.
\item \textbf{PerfChecker extensif et optimisation propre à chaque
technique.} Balayer les dimensions entières accessibles et croiser
sorties scalaires, multisorties et produit complet, métriques dégénérées
ou non orthogonales, supports, densités, horizons, types numériques,
threads et processus. Conserver les épisodes complets, la préparation,
les échantillons, la médiane et la queue, les allocations, le temps de
ramasse-miettes (GC), la taille retenue, le pic RSS, la première
compilation et les refus de budget. Examiner la stabilité du typage et
les graphes de flammes (\emph{flame graphs}) pour cibler les coûts réels,
puis mesurer de nouveau après chaque optimisation. Les résultats de
vitesse obtenus sous Étendue3D seront rejoués sur machine isolée.
Pour B1, profiler la copie et la validation répétées, puis comparer
B1W aux deux voies actuelles à sorties possédées identiques.
\item \textbf{Comparaison C++/Julia à stratégie identique.} Une fois
chaque voie qualifiée et optimisée, confronter des noyaux qui utilisent
la même stratégie, les mêmes supports, coefficients, sorties, règles
de précision et coûts de préparation ; enregistrer séparément
compilation et temps froid. La première comparaison packed ci-dessus
reste un pilote exploratoire.
\item \textbf{Classement par dimension.} Sur les campagnes précédentes
confirmées, classer les techniques à chaque dimension et pour chaque
famille de charges, en gardant les horizons, métriques et budgets comme
conditions du classement. Un croisement trop proche ou instable impose
des répétitions supplémentaires ; aucune dimension limite universelle
n'est déduite d'un cas ponctuel.
\end{enumerate}

Les tableaux de cet article sélectionnent des lignes pour l'exposé ; les
archives de cas, leurs configurations, empreintes et exclusions doivent
accompagner la version publiée pour permettre une réplication. Les tests actuels
prouvent la correction sur leurs familles, pas sur toute algèbre possible.
Julia 1.13.0 est la version retenue pour la validation courante.
Sous cette version, les suites complètes à un thread réussissent pour
\code{Garamon.jl} et \code{GaramonBench.jl} ; le sélecteur exact passe
406 assertions ciblées.
Ces réussites
valident les tests exécutés, sans qualifier les durées comme mesures
de performance puisque des processus Étendue3D partageaient la machine.
La branche publique Garamon.jl et l'environnement résolu du banc sont
archivés par empreinte pour chaque campagne ; leurs sources et paramètres
doivent accompagner toute publication des résultats. Une seule stratégie
packed possède une première
comparaison C++/Julia appariée et non confirmée en isolation. Ni catalogue
général de précompilation, ni validation GPU exhaustive, ni sélecteur HPO ne
sont revendiqués. Le
prototype de pruning est validé comme expérience comparative, sans admission
dans l'API de produit exact ni garantie de précision générale.

\subsection{Banc partageable et statut du prévol}
\label{sec:transfer}
Le paquet séparé \code{GaramonBench.jl} décrit les demandes avec
DrWatson.jl~\cite{drwatson} : paramètres, graine, algèbre, sorties,
horizon, ressources et identité des sources composent l'identifiant du
cas. Un cas achevé contient son oracle, les échantillons, les verdicts
instrumentés et les empreintes des fichiers ; la reprise contrôle ces
éléments et n'exécute pas une seconde fois le cas intact. Un cas
interrompu n'est jamais converti en résultat valide par la seule
présence de fichiers. Chaque machine reçoit une nouvelle racine de
résultats : ses temps ne sont pas importés comme points de reprise
d'une autre machine.

Le tableau~\ref{tab:preflightcoverage} distingue le prévol minimal
utilisé avant les mesures et les corpus de correction plus larges.
L'oracle vérifie la sortie mathématique et la propriété de la sortie
quand celle-ci doit survivre à l'appel. Les refus de dimension, de
mémoire ou de domaine algébrique sont aussi des résultats du prévol.
Aucune ligne de ce tableau ne classe les stratégies par vitesse.

\begin{table}[ht]
\centering\small
\begin{tabularx}{\linewidth}{@{}p{34mm}p{25mm}X@{}}
\toprule
Corpus & Cas validés & Sortie et contrôle principal \\
\midrule
Registre minimal Julia & 32/32 & Un cas par voie exécutable : 30 techniques, deux combinaisons ; 30 voies à un thread, une à quatre threads, une GPU. \\
Packed C++/Julia & 2/2 & Même algèbre euclidienne 3D et supports de grades mixtes, mêmes coefficients et matrice possédée ; oracle entier indépendant et audit des trois sources. \\
Vecteurs Julia & 312/312 & Produit possédé pour quatre chemins, treize dimensions de 2 à 128, trois signatures et deux horizons. \\
Produit triple ciblé & 420/420 & Coefficients de $(ab)c$ demandés ; comparaison avec une multiplication de mots indépendante. \\
Produit binaire ciblé et grades & 1\,176/1\,176 et 1\,344/1\,344 & Coefficient par XOR, puis produit homogène direct ou préparé ; refus avant énumération excessive. \\
Cache et indexation & 1\,116/1\,116 et 56/56 & Trace à coefficients changeants avec évictions, puis rang lexicographique et décodage des lames. \\
Stockage dense ou creux & 264/264 & Même produit complet jusqu'à 12D ; tous les coefficients comparés à l'oracle. \\
Diagnostic CPU/GPU & 8/8 & Quatre identités GPU et leurs témoins CPU ; douze traces CUDA séparent noyau, retour hôte et transferts. \\
\bottomrule
\end{tabularx}
\caption{Couverture de correction avant le benchmark isolé. Les corpus
diffèrent par leurs opérations et se chevauchent ; leurs comptes ne
s'additionnent pas pour former une mesure de performance.}
\label{tab:preflightcoverage}
\end{table}

Les 32 voies minimales ont également chacune une archive PerfChecker
complète sur la même identité de source. Le profil de la voie par
processus a nécessité une limite déclarée de 300 secondes par
collecteur ; deux tentatives à 120 secondes avaient expiré sans
fournir de pile CPU exploitable. Sur 31 voies analysées par JET,
huit ne produisent aucun signalement. Les autres totalisent
1\,290 signalements de dispatch : 23 localisés dans le banc, 262 dans
Garamon.jl et 1\,005 dans Julia ou ses dépendances. Ces comptes
décrivent des chemins de l'arbre d'inférence, avec possibles
répétitions d'une même cause ; ils ne sont ni des défauts distincts
ni des pourcentages de temps perdu. Les sorties des cas de produit
par grade, des deux produits ciblés et du sélecteur sont inférées
comme \code{Matrix\{Float64\}}, bien que des appels internes
demeurent dynamiques. La correction d'un appel interne doit donc
être choisie à partir de son profil et mesurée de nouveau, et non
du seul nombre d'alertes.

Quelques correctifs localisés illustrent cette méthode. Dans le
cache de plans, une traversée répétée de \code{summarysize} et un
tri de la liste LRU ont été réduits ; la clé inclut désormais le
type de la métrique, car confondre un rationnel et un \code{Float64}
pouvait changer le type du résultat exact. Le banc des blocs de
grades emploie un masque machine en petite dimension au lieu de
construire inutilement un \code{BigInt}. Un parcours ciblé des bits,
puis une barrière de typage, ont réduit les allocations des
coefficients binaires. Les versors factorisés évitent deux vecteurs
temporaires pour les petites métriques diagonales ; les autres
métriques gardent le chemin général. Aucune de ces corrections ne
supprime une contribution à un produit exact. Les diminutions
d'allocations ne sont pas partout accompagnées d'un gain de temps
dans les premiers échantillons, d'où la nécessité du rejeu isolé.

Le sélecteur exact fournit un autre exemple. Son estimation
conservative des coûts utilisait \code{BigInt} même à 4D. Une borne
arithmétique permet désormais le type entier machine lorsque
$n\leq12$, que l'horizon et le nombre de demandes valent au plus
4\,096, que la taille de l'algèbre est au plus 1~Gio et que celle
d'un coefficient est au plus 64~Kio sur machine 64 bits ; les cas
plus grands gardent \code{BigInt}. Les décisions et coûts des deux
branches ont été comparés dans 406 assertions ciblées. Sur le
même callback 4D, un essai PerfChecker A--B--A de 31 observations
par état donne 71,418 / 58,842 / 71,171~\(\mu\)s et
33\,064 / 24\,752 / 33\,064 octets par appel. La baisse d'allocation
est locale et le gain temporel apparent, environ 17,5~\%, reste
exploratoire ; un autre jeu de trois observations ne suffit pas
à l'étendre à d'autres scénarios. Le nombre JET de cette voie
augmente de 361 à 375 malgré la baisse d'allocation, ce qui
illustre la différence entre diagnostic statique et coût mesuré.

Pour transporter ce banc, la source C++ est figée à la révision
\code{6267161} du dépôt public~\cite{garamoncpp} dans un artefact
\code{Pkg.Artifacts}. Clang, CMake, Ninja et Eigen sont apportés
par les dépendances binaires du paquet ; les bibliothèques
générées sont construites puis nettoyées dans un espace temporaire
pour chaque dimension. Une branche publique de Garamon.jl fixe désormais
le code testé ; le banc conserve son environnement résolu et l'identité
exacte des sources dans chaque campagne. Une
reconstruction locale dans deux répertoires temporaires a passé
un prévol Julia avec contrôle des bornes, puis le prévol packed
C++/Julia 2/2 et son audit d'intégrité. Elle contrôle le transfert
des sources sur le même poste, pas encore l'installation, le pilote
GPU ni la stabilité des temps sur la machine dédiée.

Les sources de GATL, Versor, GAL, Klein, gafro et Grassmann.jl
sont également épinglées comme artefacts optionnels
\cite{gatlrepo,versorrepo,galrepo,kleinrepo,gafrorepo,grassmannrepo}.
GAL et Versor possèdent un premier adaptateur de produit de vecteurs EGA3,
avec contrôle de tous les coefficients par oracle indépendant. Leurs
stratégies internes diffèrent de celle de Garamon.jl ; ce premier graphe
décrit donc trois bibliothèques, sans isoler l'effet du langage. Les autres
comparaisons exigeront une convention
de base, une métrique, une précision, une sortie et un oracle
communs. Après transfert, la machine dédiée instanciera les
environnements Julia~1.13, exécutera les prévols CPU, quatre
threads, GPU et C++, puis lancera les grilles de mesure reprenables.
Les optimisations restantes, le choix appris, la HPO et
l'attribution SHAP attendent ces mesures isolées.

\IfFileExists{ega3_vector_libraries.pdf}{
\begin{figure}[ht]
\centering
\includegraphics[width=.88\linewidth]{ega3_vector_libraries.pdf}
\caption{Produit géométrique de 1\,024 paires de vecteurs en EGA3 : médianes
chaudes de Garamon.jl, GAL et Versor. Les stratégies internes diffèrent ; les
mesures sont décrites par leur condition dans l'archive DrWatson.}
\label{fig:external-ga-vectors}
\end{figure}
}{
\figureattente{Bibliothèques GA : vecteurs EGA3}
{Abscisse : Garamon.jl, GAL et Versor ; ordonnée : médiane chaude pour
1\,024 produits de paires identiques, avec dispersion et condition de mesure.}
{Le graphique vectoriel Makie remplacera cet emplacement après
\code{bench()} sur la machine dédiée.}{fig:external-ga-vectors}
}

\subsection{Mémoire de travail bornée : piste RAM-free}
Nous ajoutons au protocole une voie exacte à mémoire de travail réduite,
sans paramètre de taille de RAM fourni par l'utilisateur. Le nom
\emph{RAM-free} est ici une étiquette de travail : les algorithmes
\emph{cache-oblivious} optimisent les transferts sans connaître les
paramètres des caches, mais ne garantissent pas qu'une sortie complète
tienne en mémoire~\cite{frigo2012}. En effet, matérialiser un résultat
dense en dimension $n$ requiert au moins $2^n$ coefficients ; une sortie
creuse possédée requiert un espace proportionnel à son support. L'essai
doit donc comparer sortie ciblée, flux exact par blocs et représentation
comprimée, avec refus explicite si une sortie possédée excède la capacité.
Un parcours à checkpoints et recomputation, inspiré seulement dans son
compromis temps/espace des piles comprimées~\cite{baffier2018stack},
sera confronté à la récursion directe, aux workspaces et aux autres
représentations exactes. Chaque contribution et chaque signe resteront
contrôlés par oracle ; le pic RSS, les allocations, les défauts de page
et le surcoût des recomputations seront mesurés. Aucun résultat de vitesse
n'est encore attribué à cette piste.
\figureattente{RAM-free : temps contre mémoire maximale}
{Abscisse : pic de mémoire, incluant la sortie ; ordonnée : temps total,
par dimension et forme de sortie, avec refus de budget visibles.}
{Comparaison de récursion directe, flux par blocs, recomputation et
workspace lorsque les sorties sont mathématiquement identiques.}
{fig:ramfree-budget}

\FloatBarrier
\Needspace{28\baselineskip}
\section{Conclusion}
La dimension ambiante ne suffit pas à choisir un algorithme de produit
géométrique. L'audit de Garamon C++ montre la force d'un générateur
spécialisé, avec blocs de grades denses, transformations de base et
parcours récursif virtuel. La réimplémentation Julia permet d'exposer
explicitement les autres variables de décision : supports actifs,
grades, nombre de coefficients demandés, stabilité des entrées,
horizon de réutilisation et propriété des sorties. Une méthode
n'entre dans le portefeuille exact que si son domaine et toutes ses
contributions sont contrôlés.

Les premières mesures dessinent des régions plutôt qu'un classement
unique. Une sortie ciblée peut éviter le produit complet dans des
supports dispersés, mais plusieurs sorties inversent parfois le
classement entre jointure et récursion. Un plan ternaire préparé paie
un coût initial et devient favorable dans les cellules répétées de
son premier corpus ; les workspaces réduisent les allocations chaudes
lorsque la sortie peut être consommée sans copie. Les masques
\code{UInt128} changent fortement le coût de certains noyaux entre
65 et 128 dimensions, sans résoudre l'explosion d'une sortie dense.
Rang binaire, radical nilpotent, identités courtes, Pfaffien et
représentations comprimées demandent chacun un certificat de structure
et un témoin de même sortie. Les variantes approchées peuvent gagner
du temps en acceptant une erreur, mais la correction exacte par
résidu ou replay a son propre coût et ne doit jamais être choisie
silencieusement pour un produit exact.

Le prévol reproductible et les profils locaux sont complets pour
32 voies Julia et deux cas packed C++/Julia appariés. Ils valident
des contrats et rendent les coûts observables ; ils ne prouvent
pas une supériorité de Julia ou de C++, ni un seuil optimal.
Les temps de ce travail ont souvent coexisté avec d'autres calculs.
La suite expérimentale consiste à transférer les sources et
environnements figés, rejouer les oracles sur la machine dédiée,
profiler et optimiser les voies encore coûteuses, puis répéter les
grilles appariées à plusieurs dimensions, horizons et ressources.
Un sélecteur appris, son réglage d'hyperparamètres et ses explications
par SHAP ne seront justifiés qu'après ces mesures isolées.

\begin{thebibliography}{99}
\small
\bibitem{breuils2019} S. Breuils, V. Nozick et L. Fuchs, « Garamon: A Geometric Algebra Library Generator », \emph{Advances in Applied Clifford Algebras} 29, 69 (2019). \url{https://doi.org/10.1007/s00006-019-0987-7}.
\bibitem{breuils2020} S. Breuils, V. Nozick et A. Sugimoto, « Computational Aspects of Geometric Algebra Products of Two Homogeneous Multivectors » (2020). \url{https://arxiv.org/abs/2002.11313}.
\bibitem{garamoncpp} Dépôt Garamon C++, \href{https://github.com/vincentnozick/garamon/tree/6267161e35be57873fa553f266ef04f2ba32d2e5}{révision 6267161}.
\bibitem{gatlrepo} Dépôt GATL, \href{https://github.com/laffernandes/gatl/tree/1f487a70f9c0d29528b7edccf441092bdd2d34f0}{révision 1f487a7}.
\bibitem{versorrepo} Dépôt Versor, \href{https://github.com/wolftype/versor/tree/8262fce98c99ea6d9f0984d2ac17cde9fe5d0519}{révision 8262fce}.
\bibitem{galrepo} Dépôt GAL, \href{https://github.com/jeremyong/gal/tree/e859994aeb826acaf912b3efe5813582b0648291}{révision e859994}.
\bibitem{kleinrepo} Dépôt Klein, \href{https://github.com/jeremyong/klein/tree/21823fc90d191c68a4f31e0b233c58ea28ca3d5a}{révision 21823fc}.
\bibitem{gafrorepo} Dépôt gafro, \href{https://github.com/idiap/gafro/tree/5a8a5ec2cfd0d80b74017565f50824b9571a7e69}{révision 5a8a5ec}.
\bibitem{grassmannrepo} Dépôt Grassmann.jl, \href{https://github.com/chakravala/Grassmann.jl/tree/4f79a7fdb2d569bf6c13751fad3c078b6c135283}{révision 4f79a7f}.
\bibitem{cppmain} Garamon C++, \code{src/main.cpp}, chaîne de génération, \href{https://github.com/vincentnozick/garamon/blob/6267161e35be57873fa553f266ef04f2ba32d2e5/src/main.cpp}{révision 6267161}.
\bibitem{cppindex} Garamon C++, \code{src/ProductTools.cpp}, indexation des lames, \href{https://github.com/vincentnozick/garamon/blob/6267161e35be57873fa553f266ef04f2ba32d2e5/src/ProductTools.cpp}{révision 6267161}.
\bibitem{cppdispatch} Garamon C++, \code{src/ProductToString.cpp}, émission et sélection des produits, \href{https://github.com/vincentnozick/garamon/blob/6267161e35be57873fa553f266ef04f2ba32d2e5/src/ProductToString.cpp}{révision 6267161}.
\bibitem{cppmvec} Garamon C++, \code{data/Mvec.hpp}, stockage et dispatch, \href{https://github.com/vincentnozick/garamon/blob/6267161e35be57873fa553f266ef04f2ba32d2e5/data/Mvec.hpp}{révision 6267161}.
\bibitem{cppgeometric} Garamon C++, \code{data/Geometric.hpp}, produit récursif, \href{https://github.com/vincentnozick/garamon/blob/6267161e35be57873fa553f266ef04f2ba32d2e5/data/Geometric.hpp}{révision 6267161}.
\bibitem{cppmetric} Garamon C++, \code{src/MetricTools.cpp}, transformations de base par grade, \href{https://github.com/vincentnozick/garamon/blob/6267161e35be57873fa553f266ef04f2ba32d2e5/src/MetricTools.cpp}{révision 6267161}.
\bibitem{leapfrog} T. L. Veldhuizen, « Leapfrog Triejoin: a Worst-Case Optimal Join Algorithm » (2012). \url{https://arxiv.org/abs/1210.0481}.
\bibitem{sparseworkspace} G. Zhang, O. Hsu et F. Kjolstad, « Compilation of Modular and General Sparse Workspaces », \emph{Proceedings of the ACM on Programming Languages} 8 (PLDI), article 196 (2024). \url{https://doi.org/10.1145/3656426}.
\bibitem{frigo2012} M. Frigo, C. E. Leiserson, H. Prokop et S. Ramachandran, « Cache-Oblivious Algorithms », \emph{ACM Transactions on Algorithms} 8(1), article 4 (2012). \url{https://doi.org/10.1145/2071379.2071383}.
\bibitem{baffier2018stack} J.-F. Baffier, Y. Diez et M. Korman, « Experimental Study of Compressed Stack Algorithms in Limited Memory Environments », \emph{LIPIcs SEA} 103, article 19 (2018). \url{https://doi.org/10.4230/LIPIcs.SEA.2018.19}.
\bibitem{taco} F. Kjolstad, S. Kamil, S. Chou, D. Lugato et S. Amarasinghe, « The Tensor Algebra Compiler », \emph{Proceedings of the ACM on Programming Languages} 1 (OOPSLA), article 77 (2017). \url{https://doi.org/10.1145/3133901}.
\bibitem{finch} W. Ahrens, T. Fields Collin, R. Patel, K. Deeds, C. Hong et S. Amarasinghe, « Finch: Sparse and Structured Tensor Programming with Control Flow » (2024). \url{https://arxiv.org/abs/2404.16730}.
\bibitem{galley} K. Deeds, W. Ahrens, M. Balazinska et D. Suciu, « Galley: Modern Query Optimization for Sparse Tensor Programs », prépublication (2024). \url{https://arxiv.org/abs/2408.14706}.
\bibitem{gaalop} C. Schwinn, D. Hildenbrand, F. Stock et A. Koch, « Gaalop 2.0 --- a geometric algebra algorithm compiler », \emph{GraVisMa} (2010). \url{https://publica.fraunhofer.de/entities/publication/25f7307c-0502-402d-b63a-6f855b9226d8}.
\bibitem{oseledets} I. V. Oseledets, « Tensor-Train Decomposition », \emph{SIAM Journal on Scientific Computing} 33(5), 2295--2317 (2011). \url{https://doi.org/10.1137/090752286}.
\bibitem{minato} S. Minato, « Zero-Suppressed BDDs for Set Manipulation in Combinatorial Problems », \emph{DAC}, 272--277 (1993). \url{https://doi.org/10.1145/157485.164890}.
\bibitem{tinylfu} G. Einziger, R. Friedman et B. Manes, « TinyLFU: A Highly Efficient Cache Admission Policy » (2015). \url{https://arxiv.org/abs/1512.00727}.
\bibitem{sieve} Y. Zhang, J. Yang, Y. Yue, Y. Vigfusson et K. V. Rashmi, « SIEVE is Simpler than LRU: an Efficient Turn-Key Eviction Algorithm for Web Caches », \emph{NSDI}, 1229--1246 (2024). \url{https://www.usenix.org/conference/nsdi24/presentation/zhang-yazhuo}.
\bibitem{satzilla} L. Xu, F. Hutter, H. Hoos et K. Leyton-Brown, « SATzilla: Portfolio-Based Algorithm Selection for SAT », \emph{Journal of Artificial Intelligence Research} 32, 565--606 (2008). \url{https://doi.org/10.1613/JAIR.2490}.
\bibitem{aslib} B. Bischl et al., « ASlib: A Benchmark Library for Algorithm Selection » (2015). \url{https://arxiv.org/abs/1506.02465}.
\bibitem{autofolio} M. Lindauer, F. Hutter, H. Hoos et T. Schaub, « AutoFolio: An Automatically Configured Algorithm Selector », \emph{IJCAI} (2017). \url{https://doi.org/10.24963/ijcai.2017/715}.
\bibitem{egg} M. Willsey et al., « egg: Fast and Extensible Equality Saturation » (2020). \url{https://arxiv.org/abs/2004.03082}.
\bibitem{fftw} M. Frigo et S. G. Johnson, « The Design and Implementation of FFTW3 », \emph{Proceedings of the IEEE} 93(2), 216--231 (2005). \url{https://www.fftw.org/fftw-paper-ieee.pdf}.
\bibitem{juliaprecompile} Julia, manuel « Module initialization and precompilation ». \url{https://docs.julialang.org/en/v1/manual/modules/#Module-initialization-and-precompilation}.
\bibitem{julianumbers} Julia, manuel « Integers and Floating-Point Numbers ». \url{https://docs.julialang.org/en/v1/manual/integers-and-floating-point-numbers/}.
\bibitem{juliathreads} Julia, manuel « Multi-Threading ». \url{https://docs.julialang.org/en/v1/manual/multi-threading/}.
\bibitem{cudaperf} CUDA.jl, « Performance Tips ». \url{https://cuda.juliagpu.org/dev/tutorials/performance/}.
\bibitem{shap} S. M. Lundberg et S.-I. Lee, « A Unified Approach to Interpreting Model Predictions » (2017). \url{https://arxiv.org/abs/1705.07874}.
\bibitem{rhee} C.-H. Rhee et P. W. Glynn, « Unbiased Estimation with Square Root Convergence for SDE Models » (2015). \url{https://web.stanford.edu/~glynn/papers/2015/RheeG15.html}.
\bibitem{drwatson} G. Datseris, J. Isensee, S. Pech et T. Gál, « DrWatson: the perfect sidekick for your scientific inquiries », \emph{Journal of Open Source Software} 5(54), 2673 (2020). \url{https://doi.org/10.21105/joss.02673}.
\bibitem{paulib} F. Krötz, « PauLIB: A High-Performance Library for Processing Pauli Strings » (2026). \url{https://arxiv.org/abs/2605.25974}.
\bibitem{bravyitaper} S. Bravyi, J. M. Gambetta, A. Mezzacapo et K. Temme, « Tapering off qubits to simulate fermionic Hamiltonians » (2017). \url{https://arxiv.org/abs/1701.08213}.
\bibitem{ablamowicz} R. Abłamowicz, « Structure of spin groups associated with degenerate Clifford algebras », \emph{Journal of Mathematical Physics} 27, 1--6 (1986). \url{https://doi.org/10.1063/1.527361}.
\bibitem{bamberg} J. Bamberg et J. Saunders, « Exploiting degeneracy in projective geometric algebra » (2024). \url{https://arxiv.org/abs/2408.13441}.
\bibitem{prodanov} D. Prodanov, « Computation of Minimal Polynomials and Multivector Inverses in Non-Degenerate Clifford Algebras », \emph{Mathematics} 13(7), 1106 (2025). \url{https://doi.org/10.3390/math13071106}.
\bibitem{wilmotpf} G. P. Wilmot, « Clifford associativity and the Pfaffian expansion », \emph{Journal of Mathematical Physics} 29, 2346--2350 (1988). \url{https://doi.org/10.1063/1.528118}.
\bibitem{wilmot2026} G. P. Wilmot, « Construction of Exceptional Lie Algebra G2 and Non-associative Algebras Using Clifford Algebra », \emph{Advances in Applied Clifford Algebras} 36, 24 (2026). \url{https://doi.org/10.1007/s00006-025-01423-5}.
\bibitem{wimmer} M. Wimmer, « Algorithm 923: Efficient Numerical Computation of the Pfaffian for Dense and Banded Skew-Symmetric Matrices », \emph{ACM Transactions on Mathematical Software} 38(4), article 30 (2012). \url{https://doi.org/10.1145/2331130.2331138}.
\bibitem{mandolesi} A. L. G. Mandolesi, « Blade Products and the Angle Bivector of Subspaces » (2019). \url{https://arxiv.org/abs/1910.07327}.
\bibitem{bravyigaussian} S. Bravyi, « Lagrangian representation for fermionic linear optics », \emph{Quantum Information \& Computation} 5(3), 216--238 (2005). \url{https://arxiv.org/abs/quant-ph/0404180}.
\bibitem{dias} B. Dias et R. Koenig, « Classical simulation of non-Gaussian fermionic circuits », \emph{Quantum} 8, 1350 (2024). \url{https://doi.org/10.22331/q-2024-05-21-1350}.
\bibitem{ffsim} ffsim, « Orbital rotations and Givens decomposition », documentation du projet. \url{https://qiskit-community.github.io/ffsim/explanations/orbital-rotation.html}.
\bibitem{wlodarczyk} M. Włodarczyk, « Clifford Algebras Meet Tree Decompositions », \emph{Algorithmica} (2019). \url{https://doi.org/10.1007/s00453-018-0489-3}.
\bibitem{alphatensor} A. Fawzi et al., « Discovering faster matrix multiplication algorithms with reinforcement learning », \emph{Nature} 610, 47--53 (2022). \url{https://doi.org/10.1038/s41586-022-05172-4}.
\bibitem{minotaur} Z. Liu, S. Mada et J. Regehr, « Minotaur: A SIMD-Oriented Synthesizing Superoptimizer », \emph{OOPSLA} (2024). \url{https://users.cs.utah.edu/~regehr/minotaur-oopsla24.pdf}.
\bibitem{alphaevolve} « AlphaEvolve: A coding agent for scientific and algorithmic discovery » (2025). \url{https://arxiv.org/abs/2506.13131}.
\bibitem{lainewavefront} S. Laine et al., « Megakernels Considered Harmful: Wavefront Path Tracing on GPUs », \emph{High-Performance Graphics} (2013). \url{https://research.nvidia.com/sites/default/files/pubs/2013-07_Megakernels-Considered-Harmful/laine2013hpg_paper.pdf}.
\bibitem{art} V. Leis, A. Kemper et T. Neumann, « The Adaptive Radix Tree: ARTful Indexing for Main-Memory Databases », \emph{ICDE} (2013). \url{https://doi.org/10.1109/ICDE.2013.6544812}.
\bibitem{flintcrt} FLINT, « fmpz.h -- integers : Chinese remaindering », documentation du projet. \url{https://flintlib.org/doc/fmpz.html}.
\bibitem{cgalfiltered} CGAL, « \code{Filtered\_kernel} », manuel de référence. \url{https://doc.cgal.org/Manual/3.5.1/doc_html/cgal_manual/Kernel_23_ref/Class_Filtered_kernel.html}.
\end{thebibliography}

\end{document}
